\newpage
\chapter{多项式}
\setcounter{section}{0}
\section{一元多项式代数}
\subsection{定义}
\begin{formal}
\begin{definition}[一元多项式的定义]\label{def:一元多项式的定义}
    设数域$\mathbb{K}$,$x$为未定元,$a_i\in\mathbb{K}\left(0\leqslant i\leqslant n
    \right)$,则称如下形式表达式
    \[f\left(x\right)=
    a_nx^n+a_{n-1}x^{n-1}+\cdots+a_1x+a_0
    \]
    为$\mathbb{K}$上关于$x$的一元多项式.$\mathbb{K}$上的一元
    多项式全体记作$\mathbb{K}\left[x\right]$.
\end{definition}
\end{formal}
\begin{red}
\begin{remark}
    所谓“形式”二字,其一,只是形式写成和式,但加法尚未定义；其二,$x$是未定元,不是变元,因为未定元不是数,变元是数.
\end{remark}
\end{red}
\begin{formal}
\begin{definition}[次数的约定]\label{def:次数的约定}
    若$a_n\neq 0$,则$\deg f\left(x\right)=n.$

    特别地,对于常数多项式
    $f\left(x\right)=a\in\mathbb{K}$,
    若$a\neq 0$,则$\deg f\left(x\right)=0$,称为零次多项式；
    若$a=0$,称为零多项式,约定其次数为$-\infty.$
\end{definition}
\end{formal}
\begin{red}
\begin{remark}
    零多项式次数约定为$-\infty$.
\end{remark}
\end{red}
\begin{red}
    \begin{remark}
        注意,多项式不能出现$x^{-1}$这样的非法元素,也因此多项式的次数存在一个突变.
    \end{remark}
\end{red}
\subsection{运算}
\begin{formal}
\begin{definition}[多项式的相等]\label{def:多项式的相等}
    特别地,一个多项式等于零多项式当且仅当
    $a_n=a_{n-1}=\cdots=a_0=0$.

一般地,考虑两个非零多项式
\begin{align*}
&f\left(x\right)=a_nx^n+a_{n-1}x^{n-1}+\cdots+a_1x+a_0,a_n\neq 0\\
&g\left(x\right)=b_mx^m+b_{m-1}x^{m-1}+\cdots+b_1x+b_0,b_m\neq 0
\end{align*}
其中$m,n\geqslant 0.$则$f\left(x\right)=g\left(x\right)$当且仅当
$n=m$且$a_i=b_i\left(\forall 0\leqslant i\leqslant n=m\right)$
\end{definition}
\end{formal}
\begin{formal}
\begin{definition}[加法和数乘运算的定义]\label{def:加法和数乘运算的定义}
    一般地,考虑两个多项式
\begin{align*}
&f\left(x\right)=a_nx^n+a_{n-1}x^{n-1}+\cdots+a_1x+a_0\\
&g\left(x\right)=b_mx^m+b_{m-1}x^{m-1}+\cdots+b_1x+b_0
\end{align*}
不妨设$n\geqslant m$.定义加法
为\[
f\left(x\right)+g\left(x\right)=a_nx^n+\cdots+a_{m+1}x^{m+1}+\left(a_m+b_m\right)x^m+\cdot+\left(a_1+b_1\right)x+\left(a_0+b_0\right)
\]
$\forall k\in\mathbb{K}$,定义数乘
\[
k\cdot f\left(x\right)=ka_nx^n+ka_{n-1}x^{n-1}+\cdots+ka_1x+ka_0
\]
\end{definition}
\end{formal}
\begin{formal}
\begin{theorem}[线性空间性]\label{thm:线性空间性}
    $\mathbb{K}\left[x\right]$在加法和数乘下成为$\mathbb{K}$
    -线性空间.
\end{theorem}
\end{formal}
\begin{formal}
\begin{definition}[乘法的定义]\label{def:乘法的定义}
    特别地,任一多项式与零多项式的乘积均为零多项式.

    一般地,对于两个非零多项式
\begin{align*}
&f\left(x\right)=a_nx^n+a_{n-1}x^{n-1}+\cdots+a_1x+a_0,a_n\neq 0\\
&g\left(x\right)=b_mx^m+b_{m-1}x^{m-1}+\cdots+b_1x+b_0,b_m\neq 0
\end{align*}
定义它们的乘积为一个$n+m$次多项式
\begin{align*}
    h\left(x\right)&=f\left(x\right)\cdot g\left(x\right)\\
    &=c_{n+m}x^{n+m}+\cdots+c_1x+c_0
\end{align*}
其中$\displaystyle
c_{n+m}=a_nb_m,\cdots,c_k=\sum_{i+j=k}a_ib_j,\cdots,c_0=a_0b_0.
$
\end{definition}
\end{formal}
\begin{formal}
\begin{theorem}[乘法的性质]\label{thm:乘法的性质}
    乘法满足的性质如下：

    \begin{enumerate}[label={\textup{(\arabic*)}}]
        \item 交换律：$f\left(x\right)g\left(x\right)=
        g\left(x\right) f\left(x\right)$
        \item 结合律：$\left(f\left(x\right)g\left(x\right)\right)h\left(x\right)=
        f\left(x\right)\left(g\left(x\right)h\left(x\right)\right)$
        \item 分配律：$f\left(x\right)\left(g\left(x\right)+h\left(x\right)\right)=f\left(x\right)g\left(x\right)+f\left(x\right)h\left(x\right)$
        \item 乘法与数乘相容：$k\left(f\left(x\right)g\left(x\right)\right)=\left(kf\left(x\right)\right)g\left(x\right)=f\left(x\right)
        \left(kg\left(x\right)\right)$
        \item 乘法单位元：$1\cdot f\left(x\right)=f\left(x\right)$
    \end{enumerate}
\end{theorem}
\end{formal}
\subsection{交换代数}
\begin{green}
\begin{corollary}[交换代数]\label{cor:交换代数}
    $\mathbb{K}\left[x\right]$是$\mathbb{K}$-交换代数.
\end{corollary}
\end{green}
\begin{green}
\begin{corollary}[环]\label{cor:环}
    $\mathbb{K}\left[x\right]$也称$\mathbb{K}$上的一元多
    项式环.
\end{corollary}
\end{green}
\begin{red}
\begin{remark}
    在代数的定义中并未要求交换律,此处得到的交换代数是一个更强的定义.
\end{remark}
\end{red}
\begin{formal}
\begin{proposition}[积的次数]\label{prop:积的次数}
    若$f\left(x\right),g\left(x\right)\in \mathbb{K}\left[x\right]$,则\[
    \deg \left(f\left(x\right)g\left(x\right)\right)=\deg f\left(x\right)+\deg g\left(x\right)
    \]
\end{proposition}
\begin{Proof}
显然,对于零多项式也适用.
\end{Proof}
\end{formal}
\begin{green}
\begin{corollary}[一元多项式环的整性]\label{cor:一元多项式环的整性}
        若$f\left(x\right),g\left(x\right)\in \mathbb{K}\left[x\right]$,且$f\left(x\right)\neq 0,g\left(x\right)\neq 0$,
        则
        \[
        f\left(x\right)g\left(x\right)\neq 0
        \]
\end{corollary}
\begin{Proof}
\[
\deg \left(f\left(x\right)g\left(x\right)\right)=\deg f\left(x\right)+\deg g\left(x\right)\geqslant 0
\qedhere
\]
\end{Proof}
\end{green}
\begin{red}
\begin{remark}
    不是所有代数都有整性.但一元多项式环具有整性.
\end{remark}
\end{red}
\begin{brown}
\begin{example}
    $\mathbb{R}$上的代数$C\left[0,1\right]$不具有整性.
\end{example}
\end{brown}
\begin{brown}
\begin{example}
    $\mathbb{K}$上的代数$M_n\left(\mathbb{K}\right)$不具有整性.
\end{example}
\end{brown}
\begin{green}
\begin{corollary}[乘法消去律]\label{cor:乘法消去律}
    设$f\left(x\right)\neq 0 ,g\left(x\right),
    h\left(x\right)\in \mathbb{K}\left[x\right]$,且$f\left(x\right)g\left(x\right)=f\left(x\right)h\left(x\right)$,
    则\[
    g\left(x\right)=h\left(x\right)
    \]
\end{corollary}
\begin{Proof}
反证法加整性即得.
\end{Proof}
\end{green}
\begin{red}
\begin{remark}
    整性不成立,乘法消去律一般也不成立.
\end{remark}
\end{red}
\begin{formal}
\begin{proposition}[关于次数的若干结果]\label{prop:关于次数的若干结果}
    设$f\left(x\right)=a_nx^n+a_{n-1}x^{n-1}+\cdots+a_1x+a_0,g\left(x\right)=b_mx^m+b_{m-1}x^{m-1}+\cdots+b_1x+b_0
    \in \mathbb{K}\left[x\right],c\neq 0\in \mathbb{K}$,则
    \begin{enumerate}[label={\textup{(\arabic*)}}]
        \item $\deg\left(cf\left(x\right)\right)=\deg f\left(x\right)$
        \item $\deg\left(f\left(x\right)\pm g\left(x\right)\right)\leqslant \max\left\{
            \deg f\left(x\right),\deg g\left(x\right)
        \right\}$ ,等号有两种情况取到
        \[
        \begin{cases*}
            \deg f\left(x\right)\neq \deg g\left(x\right)\\
            \deg f\left(x\right)=\deg g\left(x\right),a_n+b_m\neq 0
        \end{cases*}
        \]
    \end{enumerate}
\end{proposition}
\begin{Proof}
$(2)$存在零多项式时显然成立.设序考虑最高次项即得.
\end{Proof}
\end{formal}
\newpage
\section{整除}
\subsection{定义}
\begin{formal}
\begin{definition}[整除]\label{def:整除}
    设$f\left(x\right),g\left(x\right)\in \mathbb{K}\left[x\right]$,
    若存在$h\left(x\right)\in \mathbb{K}\left[x\right]$,\,\textup{s.t.}
    \[
    f\left(x\right)=g\left(x\right)h\left(x\right)
    \]
    则称$g\left(x\right)$是$f\left(x\right)$的因式,
    或称$g\left(x\right)$可以整除$f\left(x\right)$,
    或称$f\left(x\right)$可以被$g\left(x\right)$整除,
    记作$g\left(x\right)|f\left(x\right)$；反之则称
    $g\left(x\right)$不能整除$f\left(x\right)$,或称
    $f\left(x\right)$不能被$g\left(x\right)$整除,记作
    $g\left(x\right)\nmid f\left(x\right)$.
\end{definition}
\end{formal}
\begin{red}
\begin{remark}
    零多项式的因式可以是任一多项式,但任一多项式的因式不能是零多项式.
\end{remark}
\end{red}
\begin{formal}
\begin{proposition}[整除的若干性质]\label{prop:整除的若干性质}
    设$f\left(x\right),g\left(x\right),h\left(x\right)\in \mathbb{K}\left[x\right]$,
    $c\neq 0\in\mathbb{K}$,则
    \begin{enumerate}[label={\textup{(\arabic*)}}]
        \item 若$f\left(x\right)\mid g\left(x\right)$,则$cf\left(x\right)\mid g\left(x\right)$
        \item $f\left(x\right)\mid f\left(x\right)$
        \item 若$f\left(x\right)\mid g\left(x\right),g\left(x\right)\mid h\left(x\right)$,则$f\left(x\right)\mid h\left(x\right)$
        \item 若$f\left(x\right)\mid g\left(x\right),f\left(x\right)\mid h\left(x\right)$,则$\forall u\left(x\right),v\left(x\right)\in \mathbb{K}\left[x\right]$,\,\textup{s.t.}
        \[
        f\left(x\right)\mid g\left(x\right)u\left(x\right)+h\left(x\right)v\left(x\right)
        \] 
        \item 若$f\left(x\right)|g\left(x\right)$且$g\left(x\right)\mid f\left(x\right)$,则$\exists k\in\mathbb{K}$,\,\textup{s.t.}\[
        f\left(x\right)=kg\left(x\right)
        \]
        此时$f\left(x\right),g\left(x\right)$称为相伴多项式,记作$f\left(x\right)\sim g\left(x\right)$
    \end{enumerate}
\end{proposition}
\begin{Proof}
\begin{enumerate}[label={\textup{(\arabic*)}}]
    \item $g\left(x\right)=f\left(x\right)h\left(x\right)=\left(cf\left(x\right)\right)\left(c^{-1}h\left(x\right)\right)$
    \item $f\left(x\right)=1\cdot f\left(x\right)$
    \item 代入即得
    \item 设$g\left(x\right)=f\left(x\right)p\left(x\right),h\left(x\right)=f\left(x\right)q\left(x\right)$,则
    \[
    g\left(x\right)u\left(x\right)+h\left(x\right)v\left(x\right)=f\left(x\right)\left(p\left(x\right)u\left(x\right)+q\left(x\right)v\left(x\right)\right)
    \]
    \item 设$g\left(x\right)=f\left(x\right)p\left(x\right),h\left(x\right)=f\left(x\right)q\left(x\right)$,则
    \[
    f\left(x\right)=f\left(x\right)p\left(x\right)q\left(x\right)
    \]
    $f\left(x\right)=0$显然成立.$f\left(x\right)\neq 0$时,两边取次数
    \[
    \deg f\left(x\right)=\deg f\left(x\right)+\deg \left(p\left(x\right)q\left(x\right)\right)
    \]
    于是$\deg p\left(x\right)=\deg q\left(x\right)=0$即二者均为非零常数多项式
    \qedhere
\end{enumerate}
\end{Proof}
\end{formal}
\begin{formal}
\begin{definition}[相伴多项式]\label{def:相伴多项式}
    若$f\left(x\right)\mid g\left(x\right)$且$g\left(x\right)\mid f\left(x\right)$,则
        此时$f\left(x\right),g\left(x\right)$称为相伴多项式,记作$f\left(x\right)\sim 
        g\left(x\right)$
\end{definition}
\end{formal}
\subsection{带余除法}
\begin{formal}
\begin{theorem}[带余除法]\label{thm:带余除法}
    设多项式$f\left(x\right),g\left(x\right)\in \mathbb{K}\left[x\right],
    g\left(x\right)\neq 0$,则必定存在唯一的$q\left(x\right),r\left(x\right)\in \mathbb{K}\left[x\right]$,\,\textup{s.t.}
    \[
    f\left(x\right)=g\left(x\right)q\left(x\right)+r\left(x\right)
    \]
    且有$\deg r\left(x\right)<\deg g\left(x\right)$.
\end{theorem}
\begin{Proof}
先证存在性.若$\deg f\left(x\right)<\deg g\left(x\right)$,则$q\left(x\right)=0,r\left(x\right)=f\left(x\right)$.

以下设$\deg f\left(x\right)\geqslant \deg g\left(x\right)\geqslant 0$,考虑对$\deg f\left(x\right)$归纳.
其一,若$\deg f\left(x\right)=\deg g\left(x\right)=0$,即$f\left(x\right)=a\neq 0,g\left(x\right)=b\neq 0$.则$q\left(x\right)=ab^{-1},r\left(x\right)=0$符合.

设$\deg f\left(x\right)<n$时结论成立,来证明$\deg f\left(x\right)=n$的情况.设\begin{align*}
    &f\left(x\right)=a_nx^n+a_{n-1}x^{n-1}+\cdots+a_1x+a_0,a_n\neq 0\\
    &g\left(x\right)=b_mx^m+b_{m-1}x^{m-1}+\cdots+b_1x+b_0,b_m\neq 0
\end{align*}
且$n\geqslant m.$设$f_1\left(x\right)=f\left(x\right)-a_nb_m^{-1}x^{n-m}g\left(x\right)$,于是
$\deg f_1\left(x\right)<n.$由归纳假设得知存在$q_1\left(x\right),r\left(x\right)$,\,\textup{s.t.}
\[
f_1\left(x\right)=g\left(x\right)q_1\left(x\right)+r\left(x\right)
\]
其中$\deg r\left(x\right)< \deg g\left(x\right).$即有\[
f\left(x\right)=g\left(x\right)\left(a_nb_m^{-1}x^{n-m}+q_1\left(x\right)\right)+r\left(x\right)
\]
则$q\left(x\right)=a_nb_m^{-1}x^{n-m}+q_1\left(x\right)$.
存在性证毕.

下证唯一性.考虑设出两个带余除法,作差即可.即设$f\left(x\right)=g\left(x\right)p\left(x\right)+t\left(x\right)=g\left(x\right)q\left(x\right)+r\left(x\right)$.
于是\[
g\left(x\right)\left(p\left(x\right)-q\left(x\right)\right)=t\left(x\right)-r\left(x\right)
\]
设$p\left(x\right)\neq q\left(x\right)$,对上式两边取次数
得$LHS = \deg g\left(x\right)+\deg \left(p\left(x\right)-q\left(x\right)\right)\geqslant \deg g\left(x\right)$.
$RHS=\deg\left(t\left(x\right)-r\left(x\right)\right)\leqslant \max\left\{\deg t\left(x\right),\deg r\left(x\right)\right\}<\deg g\left(x\right)$.矛盾.

证毕.
\end{Proof}
\end{formal}
\begin{green}
\begin{corollary}[整除与带余除法]\label{cor:整除与带余除法}
    设$f\left(x\right),g\left(x\right)\in \mathbb{K}\left[x\right]$,
    $g\left(x\right)\neq 0$,则$g\left(x\right)|f\left(x\right)$等价于
    $g\left(x\right)$除$f\left(x\right)$的余式为零.
\end{corollary}
\end{green}
\newpage
\section{最大公因式}
\subsection{定义}
\begin{formal}
\begin{definition}[公因式与公倍式]\label{def:公因式与公倍式}
    设$f\left(x\right),g\left(x\right)\in\mathbb{K}\left[x\right]$,
    若$d\left(x\right)\in\mathbb{K}\left[x\right]$\,\textup{s.t.}
    \[
    d\left(x\right)\mid f\left(x\right),d\left(x\right)\mid g\left(x\right)
    \]
    则称$d\left(x\right)$为$f\left(x\right)$与$g\left(x\right)$的一个公因式（公因子）；
    若$m\left(x\right)\in\mathbb{K}\left[x\right]$\,\textup{s.t.}
    \[
    f\left(x\right)\mid m\left(x\right),g\left(x\right)\mid m\left(x\right)
    \]
    则称$m\left(x\right)$是$f\left(x\right)$与$g\left(x\right)$的一个公倍式.
\end{definition}
\end{formal}
\begin{formal}
\begin{definition}[最大公因式与最小公倍式]
    若$d\left(x\right)$是$f\left(x\right)$与$g\left(x\right)$的公因式,且对于$f\left(x\right)$与$g\left(x\right)$的任一公因式$h\left(x\right)$,均有
    \[
    h\left(x\right)\mid d\left(x\right)
    \]
    则称$d\left(x\right)$是$f\left(x\right)$与$g\left(x\right)$的最大公因式,记作
    \[
    d\left(x\right)=\left(f\left(x\right),g\left(x\right)\right)=\mathrm{g.c.d.}\left(f\left(x\right),g\left(x\right)\right)
    \]

    若$m\left(x\right)$是$f\left(x\right)$与$g\left(x\right)$的公倍式,且对于$f\left(x\right)$与$g\left(x\right)$的任一公倍式$l\left(x\right)$,均有
    \[
    m\left(x\right)\mid l\left(x\right)
    \]
    则称$m\left(x\right)$是$f\left(x\right)$与$g\left(x\right)$的最小公倍式,记作
    \[
    m\left(x\right)=\left[f\left(x\right),g\left(x\right)\right]=\mathrm{l.c.m.}\left(f\left(x\right),g\left(x\right)\right)
    \]
\end{definition}
\end{formal}
\begin{red}
\begin{remark}
    约定两个零多项式的最大公因式为零多项式.
\end{remark}
\end{red}
下面利用Euclid辗转相除法证明两个多项式的最大公因式必定存在.
\begin{formal}
\begin{theorem}[最大公因式存在性]\label{thm:最大公因式存在性}
    设$f\left(x\right),g\left(x\right)\in\mathbb{K}\left[x\right]$,且二者不全为零,则$f\left(x\right),g\left(x\right)$的最大公因式
    $d\left(x\right)$必定存在且存在
    $u\left(x\right),v\left(x\right)\in\mathbb{K}\left[x\right]$\,\textup{s.t.}
    \[
    d\left(x\right)=f\left(x\right)u\left(x\right)+g\left(x\right)v\left(x\right)
    \]
\end{theorem}
\begin{Proof}
利用Euclid辗转相除法证明.

其一,若$f\left(x\right)=0$,则$\left(f\left(x\right),g\left(x\right)\right)=g\left(x\right)$,反之亦然.

设$f\left(x\right)\neq 0,g\left(x\right)\neq 0$.作带余除法
\[
f\left(x\right)=g\left(x\right)q_1\left(x\right)+r_1\left(x\right)
\]
其中$\deg r_1\left(x\right)<\deg g\left(x\right)$.于是\[
g\left(x\right)=r_1\left(x\right)q_2\left(x\right)+r_2\left(x0\right)
\]
其中$\deg r_2\left(x\right)<\deg r_1\left(x\right)$.如此重复即
\[
r_{s-2}\left(x\right)=r_{s-1}\left(x\right)q_s\left(x\right)+r_s\left(x\right)
\]
其中$\deg r_s\left(x\right)<\deg r_{s-1}\left(x\right)$.再做一次即得
\[
r_{s-1}\left(x\right)=r_s\left(x\right)q_{s+1}\left(x\right)+r_{s+1}\left(x\right)
\]
其中$\deg r_{s+1}\left(x\right)<\deg r_s\left(x\right)$.考虑到余式次数严格递减,
这个过程必定在某一步得到一个为零的余式,不妨设$r_{s+1}\left(x\right)=0,r_{s}\left(x\right)\neq 0
.$断言
\[
r_s\left(x\right)=\mathrm{g.c.d.}\left(f\left(x\right),g\left(x\right)\right)
\]
约定$r_{-1}\left(x\right)=f\left(x\right),r_0\left(x\right)=
g\left(x\right)$.
只需倒回验证即可.存在性证毕.

因为\[
r_s\left(x\right)=r_{s-2}\left(x\right)-r_{s-1}\left(x\right)q_s\left(x\right)
\]
且\[
r_{s-1}\left(x\right)=r_{s-3}\left(x\right)-r_{s-2}\left(x\right)q_{s-1}\left(x\right)
\]
则\[r_s\left(x\right)=
r_{s-3}\left(x\right)\left(-q_s\left(x\right)\right)+\left(1+q_{s-1}\left(x\right)q_s\left(x\right)\right)
r_{s-2}\left(x\right)
\]
迭代即得
\[
r_s\left(x\right)=r_{-1}\left(x\right)u\left(x\right)+r_0\left(x\right)v\left(x\right)
\]
证毕.
\end{Proof}
\end{formal}
\begin{red}
\begin{remark}
    反过来,若存在$u\left(x\right),v\left(x\right)$\,\textup{s.t.}
    \[
    d\left(x\right)=f\left(x\right)u\left(x\right)+g\left(x\right)v\left(x\right)
    \]
    却不能推出$d\left(x\right)=\left(f\left(x\right),g\left(x\right)\right)$.除非加上$d\left(x\right)\mid f\left(x\right)$且$d\left(x\right)\mid g\left(x\right)$的条件.
\end{remark}
\end{red}
\subsection{唯一性}
\begin{red}
\begin{remark}
    最大公因式必定存在,但并不一定唯一.
\end{remark}
\end{red}
\begin{formal}
\begin{theorem}[最大公因式不唯一]\label{thm:最大公因式不唯一}
    设$f\left(x\right),g\left(x\right)\in\mathbb{K}\left[x\right]$
    的两个最大公因式为
    $d_1\left(x\right),d_2\left(x\right)$,
    则$d_1\left(x\right)\mid
    d_2\left(x\right),d_2
    \left(x\right)\mid d_1\left(x\right)$.那么必
    定存在$c\neq 0\in\mathbb{K}$\,\textup{s.t.}
    \[d_2\left(x\right)=cd_1\left(x\right)\]
\end{theorem}
\end{formal}
\begin{formal}
\begin{definition}[首一多项式]\label{def:首一多项式}
    称首项系数为$1$的多项式为首一多项式.
\end{definition}
\end{formal}
\begin{red}
\begin{remark}
    约定,最大公因式均指首一多项式.在此约定下,最大公因式存在且唯一.
\end{remark}
\end{red}
\begin{formal}
\begin{theorem}[最小公倍式不唯一]\label{thm:最小公倍式不唯一}
    设多项式$f\left(x\right),g\left(x\right)\in\mathbb{K}\left[x\right]$
    的两个最小公倍式为
    $m_1\left(x\right),m_2\left(x\right)$,
    则$m_1\left(x\right)\mid 
    m_2\left(x\right),m_2
    \left(x\right)\mid m_1\left(x\right)$.那么必
    定存在$c\neq 0\in\mathbb{K}$\,\textup{s.t.}
    \[m_2\left(x\right)=cm_1\left(x\right)\]
\end{theorem}
\end{formal}
\begin{red}
\begin{remark}
    约定,最小公倍式均指首一多项式.在此约定下,最小公倍式存在且唯一.
    （存在性后续证明）
\end{remark}
\end{red}
\subsection{有限多个多项式}
\begin{formal}
\begin{definition}[有限多个多项式的公因式与最大公因式定义]\label{def:有限多个多项式的公因式与最大公因式定义}
    设存在有限多个多项式$f_1\left(x\right),f_2\left(x\right),\cdots,f_m\left(x\right)\in\mathbb{K}\left[x\right]\left(m\geqslant 2\right)$
,若$d\left(x\right)\mid f_i\left(x\right)\left(\forall 1\leqslant i\leqslant m\right)$,则称$d\left(x\right)$为
$f_1\left(x\right),f_2\left(x\right),\cdots,f_m\left(x\right)$的一个公因式.
进一步的,若$f_1\left(x\right),f_2\left(x\right),\cdots,f_m\left(x\right)$的任一公因式
$h\left(x\right)$均有
\[
h\left(x\right)\mid d\left(x\right)
\]
则称$d\left(x\right)$为$f_1\left(x\right),f_2\left(x\right),\cdots,f_m\left(x\right)$的最大公因式,
记作\[
d\left(x\right)=\left(
    f_1\left(x\right),f_2\left(x\right),\cdots,f_m\left(x\right)
\right)\left(m\geqslant 2\right)
\]
\end{definition}
\end{formal}
\begin{formal}
\begin{definition}[有限多个多项式的公倍式与最小公倍式定义]\label{def:有限多个多项式的公倍式与最小公倍式定义}
    设存在有限多个多项式$f_1\left(x\right),f_2\left(x\right),\cdots,f_m\left(x\right)\in\mathbb{K}\left[x\right]\left(m\geqslant 2\right)$
,若$f_i\left(x\right)|m\left(x\right)\left(\forall 1\leqslant i\leqslant m\right)$,则称$m\left(x\right)$为
$f_1\left(x\right),f_2\left(x\right),\cdots,f_m\left(x\right)$的一个公倍式.
进一步的,若$f_1\left(x\right),f_2\left(x\right),\cdots,f_m\left(x\right)$的任一公倍式
$l\left(x\right)$均有
\[
m\left(x\right)\mid l\left(x\right)
\]
则称$m\left(x\right)$为
$f_1\left(x\right),f_2\left(x\right),\cdots,f_m\left(x\right)$的最小公倍式,
记作\[
m\left(x\right)=\left[
    f_1\left(x\right),f_2\left(x\right),\cdots,f_m\left(x\right)
\right]\left(m\geqslant 2\right)
\]
\end{definition}
\end{formal}
\begin{green}
\begin{lemma}[多个多项式的最大公因式求取]\label{lem:多个多项式的最大公因式求取}
    设多项式$f_1\left(x\right),f_2\left(x\right),\cdots,f_m\left(x\right)\in\mathbb{K}\left[x\right]\left(m\geqslant 3\right)$,
    则\[
    \left(f_1\left(x\right),f_2\left(x\right),\cdots,f_m\left(x\right)\right)=\left(\left(f_1\left(x\right),f_2\left(x\right)\right),f_3\left(x\right),\cdots,f_m\left(x\right)\right)
    \]
\end{lemma}
\begin{Proof}
设$d\left(x\right)=\left(f_1\left(x\right),f_2\left(x\right),\cdots,f_m\left(x\right)\right),d_{12}\left(x\right)=\left(f_1\left(x\right),f_2\left(x\right)\right)$.
只需证\[
d\left(x\right)=\left(d_{12}\left(x\right),f_3\left(x\right),\cdots,f_m\left(x\right)\right)
\]
一方面,$d\left(x\right)|f_i\left(x\right)\left(\forall 1\leqslant i\leqslant m\right)$,
故$d\left(x\right)|d_{12}\left(x\right),d\left(x\right)|f_i\left(x\right)\left(\forall 3\leqslant i\leqslant m\right)$.另一方面,
任取$h\left(x\right)$为$d_{12}\left(x\right),f_3\left(x\right),\cdots,f_m\left(x\right)$,则
$h\left(x\right)\mid d_{12}\left(x\right)$,故$h\left(x\right)\mid f_1\left(x\right)$且
$h\left(x\right)\mid f_2\left(x\right)$,于是$h\left(x\right)\mid f_i\left(x\right)\left(\forall 1\leqslant i\leqslant m\right)$,因此$h\left(x\right)\mid d\left(x\right)$.证毕.
\end{Proof}
\end{green}
\subsection{互素}
\begin{formal}
\begin{definition}[互素的定义]\label{def:互素的定义}
    设$f\left(x\right),g\left(x\right)\in\mathbb{K}\left[x\right]$,若$\mathrm{g.c.d.}\left(f\left(x\right),g\left(x\right)\right)=1$,称$f\left(x\right),g\left(x\right)$互素.
\end{definition}
\end{formal}
\begin{formal}
\begin{theorem}[互素判定]\label{thm:互素判定}
    设$f\left(x\right),g\left(x\right)\in\mathbb{K}\left[x\right]$,
    则$f\left(x\right)$与$g\left(x\right)$互素等价于
    存在$u\left(x\right),v\left(x\right)\in\mathbb{K}\left[x\right]$\,\textup{s.t.}
    \[
    f\left(x\right)u\left(x\right)+g\left(x\right)v\left(x\right)=1
    \]
\end{theorem}
\begin{Proof}
必要性,即已知互素,$d\left(x\right)=1$,由前定理得证.

充分性,设$d\left(x\right)=\left(f\left(x\right),g\left(x\right)\right)$,则
$d\left(x\right)\mid f\left(x\right)u\left(x\right)+g\left(x\right)v\left(x\right)=1$,即$\exists
t\left(x\right)\in\mathbb{K}\left[x\right]$\textup{s.t.}$1=d\left(x\right)t\left(x\right)$.两边同时取次数即得
$d\left(x\right)=a\neq 0$,考虑到首一多项式,$d\left(x\right)=1$.

证毕.
\end{Proof}
\end{formal}
\begin{green}
\begin{corollary}[素因子]\label{cor:素因子}
    设$f_1\left(x\right)\mid f\left(x\right)$且$f_2\left(x\right)\mid f\left(x\right)$,若$\left(f_1\left(x\right),f_2\left(x\right)\right)=1$,则\[
    f_1\left(x\right)f_2\left(x\right)\mid f\left(x\right)
    \]
\end{corollary}
\begin{Proof}
设$u\left(x\right),v\left(x\right)\in\mathbb{K}\left[x\right]$\textup{s.t.}
\[
f_1\left(x\right)u\left(x\right)+f_2\left(x\right)v\left(x\right)=1
\]
即\[
f_1\left(x\right)u\left(x\right)f\left(x\right)+f_2\left(x\right)v\left(x\right)f\left(x\right)=f\left(x\right)
\]
明显$f_1\left(x\right)f_2\left(x\right)\mid f_1\left(x\right)u\left(x\right)f\left(x\right)$,$f_1\left(x\right)f_2\left(x\right)\mid
f_2\left(x\right)v\left(x\right)f\left(x\right)$.证毕.
\end{Proof}
\end{green}
\begin{green}
\begin{corollary}[分解]\label{cor:分解}
    若$\left(f\left(x\right),g\left(x\right)\right)=1$,且$f\left(x\right)\mid g\left(x\right)h\left(x\right)$,则\[
    f\left(x\right)\mid h\left(x\right)
    \]
\end{corollary}
\begin{Proof}
类似地
\[
f\left(x\right)u\left(x\right)h\left(x\right)+g\left(x\right)v\left(x\right)h\left(x\right)=h\left(x\right)
\qedhere
\]
\end{Proof}
\end{green}
\begin{green}
\begin{corollary}[素因子之积]\label{cor:素因子之积}
    若$\left(f_1\left(x\right),g\left(x\right)\right)=1,\left(f_2\left(x\right),g\left(x\right)\right)=1$,则
    \[
    \left(f_1\left(x\right)f_2\left(x\right),g\left(x\right)\right)=1
    \]
\end{corollary}
\begin{Proof}
因为\[
f_1\left(x\right)u_1\left(x\right)+g\left(x\right)v_1\left(x\right)=1
\]
\[
f_2\left(x\right)u_2\left(x\right)+g\left(x\right)v_2\left(x\right)=1
\]
于是\begin{align*}
    &\left(f_1\left(x\right)f_2\left(x\right)\right)\left(u_1\left(x\right)u_2\left(x\right)\right)\\
    +&g\left(x\right)\left[f_1\left(x\right)u_1\left(x\right)v_2\left(x\right)+f_2\left(x\right)u_2\left(x\right)v_1\left(x\right)+g\left(x\right)v_1\left(x\right)v_2\left(x\right)\right]=1
\qedhere
\end{align*}
\end{Proof}
\end{green}
\begin{green}
\begin{corollary}[最大公因式]\label{cor:最大公因式}
    设$d\left(x\right)=\left(f\left(x\right),g\left(x\right)\right)$,且$f\left(x\right)=
    f_1\left(x\right)d\left(x\right),g\left(x\right)=g_1\left(x\right)d\left(x\right)$,则\[
    \left(f_1\left(x\right),g_1\left(x\right)\right)=1
    \]
\end{corollary}
\begin{Proof}
辗转相除代入即可.
\end{Proof}
\end{green}
\begin{green}
\begin{corollary}[外乘因子]\label{cor:外乘因子}
    设$\left(f\left(x\right),g\left(x\right)\right)=d\left(x\right)$,则\[
    \left(t\left(x\right)f\left(x\right),t\left(x\right)g\left(x\right)\right)=t\left(x\right)d\left(x\right)
    \]
\end{corollary}
\begin{Proof}
因为\[
f\left(x\right)u\left(x\right)+g\left(x\right)v\left(x\right)=d\left(x\right)
\]
则\[
f\left(x\right)t\left(x\right)u\left(x\right)+g\left(x\right)t\left(x\right)v\left(x\right)=t\left(x\right)d\left(x\right)
\]
任取$f\left(x\right)t\left(x\right),g\left(x\right)t\left(x\right)$公因子为$h\left(x\right)$,则
\[
h\left(x\right)\mid t\left(x\right)d\left(x\right)
\qedhere
\]
\end{Proof}
\end{green}
\subsection{最小公倍式}
\begin{formal}
\begin{theorem}[最小公倍式存在性]\label{thm:最小公倍式存在性}
    设非零多项式$f\left(x\right),g\left(x\right)$,则\[
    f\left(x\right)g\left(x\right)\sim \left(f\left(x\right)g\left(x\right)\right)\left[f\left(x\right),g\left(x\right)\right]
    \]
\end{theorem}
\begin{Proof}
设最大公因式$d\left(x\right)=\left(f\left(x\right),g\left(x\right)\right),
f\left(x\right)=f_1\left(x\right)d\left(x\right),g\left(x\right)=g_1\left(x\right)d\left(x\right)$,于是$\left(
    f_1\left(x\right),g_1\left(x\right)
\right)=1.$于是\[
\frac{
    f\left(x\right)g\left(x\right)
}{\left(f\left(x\right),g\left(x\right)\right)}=f_1\left(x\right)d\left(x\right)g_1\left(x\right)
\]
只需证明$f_1\left(x\right)d\left(x\right)g_1\left(x\right)$为最小公倍式.

一方面,显然这是$f\left(x\right),g\left(x\right)$的一个公倍式.另一方面,
任取$l\left(x\right)$为$f\left(x\right),g\left(x\right)$的一个公倍式.
因为$l\left(x\right)=f\left(x\right)u\left(x\right)=g\left(x\right)
v\left(x\right)$.于是
\[
f_1\left(x\right)u\left(x\right)=g_1\left(x\right)v\left(x\right)
\]因为$\left(f_1\left(x\right),g_1\left(x\right)\right)=1$,故\[
g_1\left(x\right)\mid u\left(x\right)
\]
设$u\left(x\right)=g_1\left(x\right)p\left(x\right)$,于是$l\left(x\right)=f_1\left(x\right)d\left(x\right)g_1\left(x\right)p\left(x\right)$.

证毕.
\end{Proof}
\end{formal}

\subsection{Chinese Remainden Theorem}
\begin{formal}
\begin{theorem}[Chinese Remainden Theorem]\label{thm:Chinese Remainden Theorem}
中国剩余定理，又叫孙子定理.

设$g_1\left(x\right),\cdots,g_m\left(x\right)$为两两互素的多项式,
$r_1\left(x\right),\cdots,r_m\left(x\right)\in\mathbb{K}\left[x\right]$,则存在$f\left(x\right)\in\mathbb{K}\left[x\right]$使得
\[
f\left(x\right)=g_i\left(x\right)q_i\left(x\right)+r_i\left(x\right)\left(\forall 1\leqslant i\leqslant m\right)
\]
\end{theorem}
\begin{Proof}
构造\[f_i\left(x\right)=g_i\left(x\right)p_i\left(x\right)+1\left(\forall 1\leqslant i\leqslant m\right)\]
且\[
g_j\left(x\right)\mid f_i\left(x\right)\left(j\neq i\right)
\]
那么$\displaystyle
f\left(x\right)=\sum_{i=1}^{m}r_i\left(x\right)f_i\left(x\right)
.$

$f_1\left(x\right)$满足如下条件：
\begin{align*}
&    f_1\left(x\right)\equiv 1\pmod {g_1\left(x\right)}\\
 &   f_1\left(x\right)\equiv 0\pmod {g_j\left(x\right),j\neq 1}
\end{align*}
因为$\left(g_1\left(x\right),g_j\left(x\right)\right)=1\left(j\neq 1\right)$.故$\exists u_j\left(x\right),v_j\left(x\right)\in\mathbb{K}\left[x\right]$\,\textup{s.t.}
\[
g_1\left(x\right)u_j\left(x\right)+g_j\left(x\right)v_j\left(x\right)=1
\]
于是构造$f_1\left(x\right)$如下：
\begin{align*}
    f_1\left(x\right)&=g_2\left(x\right)v_2\left(x\right)\cdots g_n\left(x\right)v_n\left(x\right)\\
    &=\left(1-g_1\left(x\right)u_2\left(x\right)\right)\cdots\left(1-g_1\left(x\right)u_n\left(x\right)\right)
\end{align*}
其余同理,证毕.
\end{Proof}
\end{formal}
\newpage
\section{因式分解}
\subsection{可约与不可约}
\begin{formal}
\begin{definition}[可约多项式与不可约多项式]\label{def:可约多项式与不可约多项式}
    设$f\left(x\right)\in\mathbb{K}\left[x\right]$,
    $\deg f\left(x\right)\geqslant 1$.若存在$g\left(x\right),h\left(x\right)\in\mathbb{K}\left[x\right]$,且
    $\deg g\left(x\right)<\deg f\left(x\right),\deg h\left(x\right)<\deg f\left(x\right)$,\textup{s.t.}
    \[f\left(x\right)=g\left(x\right)h\left(x\right)\]
    那么称$f\left(x\right)$在$\mathbb{K}\left[x\right]$上是可约的.
    否则称为是不可约的.
\end{definition}
\end{formal}
\begin{red}
\begin{remark}
    不管在哪个域上,一次多项式总是不可约的.
\end{remark}
\end{red}
\begin{formal}
\begin{definition}[可约的等价定义]\label{def:可约的等价定义}
    设$f\left(x\right)\in\mathbb{K}\left[x\right]$,
    $\deg f\left(x\right)\geqslant 1$.若存在多项式$g\left(x\right),h\left(x\right)\in\mathbb{K}\left[x\right]$,且
    $\deg g\left(x\right)>0,\deg h\left(x\right)>0$,\textup{s.t.}
    \[f\left(x\right)=g\left(x\right)h\left(x\right)\]
    那么称$f\left(x\right)$在$\mathbb{K}\left[x\right]$上是可约的.
\end{definition}
\end{formal}
\begin{formal}
\begin{definition}[不可约的等价定义]\label{def:不可约的等价定义}
    设$f\left(x\right)\in\mathbb{K}\left[x\right],\deg f\left(x\right)\geqslant 1$,则$f\left(x\right)$不可约等价于
    $f\left(x\right)$的因式有且仅有$c\neq 0$和$cf\left(x\right)$.
\end{definition}
\end{formal}
\begin{red}
\begin{remark}
    多项式的可约与否取决于数域的选择.
\end{remark}
\end{red}
\begin{brown}
\begin{example}
    $f\left(x\right)=x^2-2$在$\mathbb{Q}\left[x\right]$不可约,在$\mathbb{R}\left[x\right]$可约.
\end{example}
\end{brown}
\begin{green}
\begin{lemma}[不可约多项式与其他多项式]\label{lem:不可约多项式与其他多项式}
    设$p\left(x\right)$为$\mathbb{K}\left[x\right]$
    上的一个不可约多项式,$f\left(x\right)\in \mathbb{K}\left[x\right]$,
    则要么$\left(p\left(x\right),f\left(x\right)\right)=1$,要么
    $p\left(x\right)\mid f\left(x\right)$.
\end{lemma}
\begin{Proof}
设$\left(p\left(x\right),f\left(x\right)\right)=d\left(x\right)$,则$d\left(x\right)|p\left(x\right)$.
那么显然,$d\left(x\right)=1$或者$d\left(x\right)=cp\left(x\right)$（首一多项式）.
\end{Proof}
\end{green}
\begin{formal}
\begin{theorem}[不可约多项式的素性]\label{thm:不可约多项式的素性}
    设不可约多项式$p\left(x\right)\in\mathbb{K}\left[x\right]$,
    且有多项式$f\left(x\right),g\left(x\right)\in\mathbb{K}\left[x\right]$且$p\left(x\right)\mid f\left(x\right)g\left(x\right)$,则
    $p\left(x\right)\mid f\left(x\right)$或$p\left(x\right)\mid g\left(x\right)$.
\end{theorem}
\begin{Proof}
由\cref{lem:不可约多项式与其他多项式}及互素多项式性质立得.
\end{Proof}
\end{formal}
\begin{red}
\begin{remark}
    对于首一不可约多项式,互异即是互素.
\end{remark}
\end{red}
\begin{green}
\begin{corollary}[推广到有限多个多项式]\label{cor:推广到有限多个多项式}
    设不可约多项式$p\left(x\right)$且\[
    p\left(x\right)\mid f_1\left(x\right)f_2\left(x\right)\cdots f_m\left(x\right)
    \]
    则$\exists 1\leqslant i\leqslant m$,\textup{s.t.}$p\left(x\right)\mid f_i\left(x\right)$.
\end{corollary}
\begin{Proof}
归纳立得.
\end{Proof}
\end{green}
\subsection{因式分解}
\begin{formal}
\begin{theorem}[因式分解定理]\label{thm:因式分解定理}
    设$f\left(x\right)\in\mathbb{K}\left[x\right]$且$\deg f\left(x\right)\leqslant 1$,则
    \begin{enumerate}[label={\textup{(\arabic*)}}]
        \item 存在性：$f\left(x\right)$一定可分解为$\mathbb{K}$上有限多个不可约多项式之积
        \item 唯一性：设\[
        f\left(x\right)=p_1\left(x\right)p_2\left(x\right)\cdots p_s\left(x\right)=q_1\left(x\right)q_2\left(x\right)\cdots q_t\left(x\right)
        \]
        是$f\left(x\right)$的两个不可约分解,即$p_i\left(x\right),q_j\left(x\right)$均为$\mathbb{K}$上的次数大于零的不可约多项式,则一定有$s=t$且在适当调换
        因式的次序后,有\[
        p_i\left(x\right)\sim q_i\left(x\right)\left(\forall 1\leqslant i\leqslant s=t\right)
        \]
    \end{enumerate}
\end{theorem}
\begin{Proof}\,

\begin{enumerate}[label={\textup{(\arabic*)}}]
    \item 对$\deg f\left(x\right)=n$归纳,$n=1$显然成立.设$\deg f\left(x\right)<n$时成立,下证$\deg f\left(x\right)=n$.
    $f\left(x\right)$不可约时显然成立.设$f\left(x\right)$可约.则
    $f\left(x\right)=g\left(x\right)h\left(x\right),g\left(x\right),h\left(x\right)\in\mathbb{K}\left[x\right],\deg g\left(x\right)<n,\deg h\left(x\right)<n$.
    则由归纳假设有
    $g\left(x\right)=p_1\left(x\right)\cdots p_s\left(x\right),h\left(x\right)=q_1\left(x\right)q_2\left(x\right)\cdots q_t\left(x\right)$,其中$p_i\left(x\right),q_j\left(x\right)$均为$\mathbb{K}\left[x\right]$上的不可约多项式.
    于是存在性证毕.

    \item 对不可约因式个数$s$归纳,$s=1$时$f\left(x\right)=p_1\left(x\right)=q_1\left(x\right)\cdots q_t\left(
        x\right)$,$t\geqslant 2$时$p_1\left(x\right)$可约,矛盾.于是$t = 1.$
    即$s=t$且$p_1\left(x\right)=q_1\left(x\right)$.

    设不可约因式个数小于$s$时成立,下证等于时成立.因为\[
            f\left(x\right)=p_1\left(x\right)p_2\left(x\right)\cdots p_s\left(x\right)=q_1\left(x\right)q_2\left(x\right)\cdots q_t\left(x\right)
    \]
    故$p_1\left(x\right)\mid q_1\left(x\right)q_2\left(x\right)\cdots q_t\left(x\right)$,由素性知,$\exists 1\leqslant i\leqslant t$,\textup{s.t.}
    $p_1\left(x\right)\mid q_i\left(x\right).$不妨调整为$i=1$.即$q_1\left(x\right)=p_1\left(x\right)h\left(x\right).$由$p_1\left(x\right)$的不可约性知$h\left(x\right)=c\neq 0.$
    即$q_1\left(x\right)=cp_1\left(x\right)\Longrightarrow p_1\left(x\right)\sim q_1\left(x\right)$.于是得到
    \[
    p_2\left(x\right)\cdots p_s\left(x\right)=cq_2\left(x\right)\cdots q_t\left(x\right)
    \]
    则由归纳假设有,$s-1=t-1\Longrightarrow s=t$且调换次序下$p_i\left(x\right)\sim q_i\left(x\right)\left(\forall 2\leqslant i \leqslant s=t\right)$.

    于是得证.\qedhere
    \end{enumerate}
\end{Proof}
\end{formal}
\begin{green}
\begin{corollary}[标准因式分解]\label{cor:标准因式分解}
    设$f\left(x\right)\in\mathbb{K}\left[x\right],\deg f\left(x\right)\geqslant 1$,则有如下标准因式分解
    \[
    f\left(x\right)=cp_1\left(x\right)^{e_1}p_2\left(x\right)^{e_2}\cdots p_m\left(x\right)^{e_m}
    \]
    其中$c\neq 0\in\mathbb{K}$,$p_i\left(x\right)\left(1\leqslant i\leqslant m\right)$为互异的首一不可约多项式,$\forall e_i\geqslant 1\left(1\leqslant i\leqslant m\right)$.
\end{corollary}
\end{green}
\begin{formal}
\begin{definition}[重因式]\label{def:重因式}
    若标准因式分解中$e_1>1\left(\forall 1\leqslant i\leqslant m\right)$,则$p_i\left(x\right)$称为$f\left(x\right)$的$e_i$重因式.$e_i=1$时称为单因式.
\end{definition}
\end{formal}
\begin{green}
\begin{corollary}[公共因式分解]\label{cor:公共因式分解}
    设$f\left(x\right),g\left(x
    \right)\in\mathbb{K}\left[x\right]$,则添加某些不可约多项式的零次幂后,有如下公共因式分解
    \begin{align*}
        &f\left(x\right)=c_1p_1\left(x\right)^{e_1}p_2\left(x\right)^{e_2}\cdots p_n\left(x\right)^{e_n};
        \\
        &g\left(x\right)=c_2p_1\left(x\right)^{f_1}p_2\left(x\right)^{f_2}\cdots p_n\left(x\right)^{f_n}.
    \end{align*}
    其中$c_1\neq 0,c_2\neq 0\in\mathbb{K}$,$p_i\left(x\right)\left(1\leqslant i\leqslant n\right)$为互异的首一不可约多项式,$e_i\geqslant 0,f_i\geqslant 0\left(\forall 1\leqslant i\leqslant n\right)$.
\end{corollary}
\end{green}
\begin{green}
\begin{corollary}[公共因式分解的若干结论]\label{cor:公共因式分解的若干结论}
    \,

    \begin{enumerate}[label={\textup{(\arabic*)}}]
        \item $g\left(x\right)\mid f\left(x\right)\Longleftrightarrow f_i\leqslant e_i\left(\forall 1\leqslant i\leqslant 
        n\right)$
        \item 令$k_i=\min\left\{e_i,f_i\right\},l_i=\max\left\{e_i,f_i\right\}\left(\forall 1\leqslant i\leqslant n\right)$,则有
        \begin{align*}
            &\left(f\left(x\right),g\left(x\right)\right)=p_1\left(x\right)^{k_1}p_2\left(x\right)^{k_2}\cdots p_n\left(x\right)^{k_n}\\
            &\left[f\left(x\right),g\left(x\right)\right]=p_1\left(x\right)^{l_1}p_2\left(x\right)^{l_2}\cdots p_n\left(x\right)^{l_n}
        \end{align*}
        \item 由上面的结论,非常容易证明
        \[f\left(x\right)g\left(x\right)\sim \left(f\left(x\right),g\left(x\right)\right)\left[f\left(x\right),g\left(x\right)\right]\]
    \end{enumerate}
\end{corollary}
\end{green}
\begin{red}
\begin{remark}
    考虑计算某一多项式是否具有重因式的问题,尽管我们可以计算出它的标准因式分解,但前提是找到合适的不可约多项式,这是极其困难的一件事.
\end{remark}
\end{red}
\subsection{形式导数与重因式判定}
\begin{formal}
\begin{definition}[形式导数定义]\label{def:形式导数定义}
    设$f\left(x\right)=a_nx^n+a_{n-1}x^{n-1}+\cdots+a_1x+a_0\in\mathbb{K}\left[x\right]$,则
    \[
    f'\left(x\right)=na_{n}x^{n-1}+\left(n-1\right)a_{n-1}x^{n-2}+\cdots+a_1
    \]
    称为其形式导数.
\end{definition}
\end{formal}
\begin{formal}
\begin{theorem}[形式导数的若干性质]\,

    \begin{enumerate}[label={\textup{(\arabic*)}}]
        \item $\left(f\left(x\right)+g\left(x\right)\right)'=f'\left(x\right)+g'\left(x\right)$
        \item $\left(cf\left(x\right)\right)'=cf'\left(x\right)$
        \item $\left(f\left(x\right)g\left(x\right)\right)'=f'\left(x\right)g\left(x\right)+f\left(x\right)g'\left(x\right)$
        \item $\left(f\left(x\right)^m\right)'=mf\left(x\right)^{m-1}f'\left(x\right)$
        \item 若$\deg f\left(x\right)\geqslant 1$,则$\deg f'\left(x\right)=\deg f\left(x\right)-1$,那么
        \[
        f\left(x\right)\nmid f'\left(x\right)
        \]
    \end{enumerate}
\end{theorem}
\end{formal}
\begin{green}
\begin{corollary}[相同不可约因子]\label{cor:相同不可约因子}
    设多项式$f\left(x\right)\in\mathbb{K}\left[x\right],\deg f\left(x\right)\geqslant 1$,其形式导数为$f'\left(x\right)$,$d\left(x\right)=\left(f\left(x\right),f'\left(x\right)\right)$,则$\displaystyle
    \frac{f\left(x\right)}{d\left(x\right)}$一定没有重因式且与$f\left(x\right)$有相同的不可约因子（不计重数）.
\end{corollary}
\begin{Proof}
设标准分解
\[
f\left(x\right)=cp_1\left(x\right)^{e_1}p_2\left(x\right)^{e_2}\cdots p_m\left(x\right)^{e_m}
\]
其中$c\neq 0\in\mathbb{K}$,$p_i\left(x\right)\left(1\leqslant i\leqslant m\right)$为互异的首一不可约多项式,$e_i\geqslant 1\left(1\leqslant i\leqslant m\right)$.
求导即有
\begin{align*}
    f'\left(x\right)=&ce_1p_1\left(x\right)^{e_1-1}p_2\left(x\right)^{e_2}\cdots p_m\left(x\right)^{e_m}p_1'\left(x\right)\\
    &+ce_2p_1\left(x\right)^{e_1}p_2\left(x\right)^{e_2-1}\cdots p_m\left(x\right)^{e_m}p_2'\left(x\right)\\
    &+\cdots\\
    &+ce_mp_1\left(x\right)^{e_1}p_2\left(x\right)^{e_2}\cdots p_m\left(x\right)^{e_m-1}p_m'\left(x\right)
\end{align*}
故$p_1\left(x\right)^{e_1-1}p_2\left(x\right)^{e_2-1}\cdots p_m\left(x\right)^{e_m-1}$是$f\left(x\right),f'\left(x\right)$的一个公因式.下面证明它是最大公因式.

设$d\left(x\right)=\left(f\left(x\right),f'\left(x\right)\right)=
p_1\left(x\right)^{k_1}p_2\left(x\right)^{k_2}\cdots p_m\left(x\right)^{k_m}$
.$e_i-1\leqslant k_i\leqslant 
e_i\left(\forall 1\leqslant i
\leqslant m\right)$,断言全部取下界.

考虑反证法,设$k_1=e_1$,即$p_1\left(x\right)^{e_1}\mid f'\left(x\right)$,显然矛盾.于是
\[
d\left(x\right)=p_1\left(x\right)^{e_1-1}p_2\left(x\right)^{e_2-1}\cdots p_m\left(x\right)^{e_m-1}
\]
则$\displaystyle
\frac{f\left(x\right)}{d\left(x\right)}=cp_1\left(x\right)p_2\left(x\right)
\cdots p_m\left(x\right)$无重因式且与$f\left(x\right)$有相同的不可约因子.证毕.
\end{Proof}
\end{green}
\begin{formal}
\begin{theorem}[重因式的判定定理]\label{thm:重因式的判定定理}
    设$f\left(x\right)\in\mathbb{K}\left[x\right],\deg f\left(x\right)\geqslant 1$,则$f\left(x\right)$
    无重因式等价于
    $\left(f\left(x\right),f'\left(x\right)\right)=1.$
\end{theorem}
\end{formal}
\begin{red}
\begin{remark}
    这表明,计算一个多项式是否存在重因式时,可以
    避开标准分解（实际上找到标准分解不见得简单,甚至会非常难,因为想要证明或判定
    一个多项式在域上不可约实际上异常困难,目前也仅有一些充分条件而没有必要条件
    ）,转而利用辗转相除法求得
    该多项式与其导数的公因式从而根据
    上面的定理判定.
\end{remark}
\end{red}
\newpage
\section{多项式函数}
\subsection{根}
\begin{formal}
\begin{definition}[多项式函数]\label{def:多项式函数}
    设$f\left(x\right)=a_nx^n+a_{n-1}x^{n-1}+\cdots+a_1x+a_0\in\mathbb{K}\left[x\right]$,则其是一个函数
    \begin{align*}
        f:\mathbb{K}&\longrightarrow \mathbb{K}\\
        b&\longmapsto f\left(b\right)
    \end{align*}
    $f\left(b\right):=
        a_nb^n+a_{n-1}b^{n-1}+\cdots a_1b+a_0$称为$f\left(x\right)$在$b$点的取值.

        于是对于一个多项式$f$,我们可以自然地视作一个多项式函数,记作$\overline{f}$.
\end{definition}
\end{formal}
\begin{red}
\begin{remark}
一个问题是,两个多项式作为
多项式相等与作为
多项式函数相等是否是一回事？前者推后者显然,
因为我们给出的多项式形式定义足够强.我们下面来考虑多项式相等得出多项式函数相等
这件事.
\end{remark}
\end{red}
\begin{red}
\begin{remark}
    一般地,多项式$f\in\mathbb{K}\left[x\right]$诱导的多项式函数记作$\overline{f}$,即我们需要证明
    \[
    f=g\Longleftrightarrow \overline{f}=\overline{g}
    \]
\end{remark}
\end{red}
\begin{red}
\begin{remark}
    在数域上,这一问题可能有一点“显然”,
    但实际上,数域只是一种特殊的域,
    后面我们将说明,
    这个等价关系在一般的域上不见得成立.
\end{remark}
\end{red}
\begin{formal}
\begin{definition}[根的定义]\label{def:根的定义}
    设$f\left(x\right)\neq 0\in\mathbb{K}
    \left[x\right]$,若$b\in\mathbb{K}$,\textup{s.t.}$f\left(b\right)=0$,则称
    $b$是$f\left(x\right)$的根或零点.
\end{definition}
\end{formal}
\begin{formal}
\begin{theorem}[余数定理]\label{thm:余数定理}
    设$f\left(x\right)\in\mathbb{K}\left[x\right],b\in\mathbb{K}$,则存在$g\left(x\right)\in\mathbb{K}$,\textup{s.t.}
    \[
    f\left(x\right)=\left(x-b\right)g\left(x\right)+f\left(b\right)
    \]
    特别地,$b$是$f\left(x\right)$的根当且仅当$\left(x-b\right)\mid f\left(x\right)$.
\end{theorem}
\begin{Proof}
首先有带余除法
\[
f\left(x\right)=\left(x-b\right)q\left(x\right)+r\left(x\right)
\]
其中$\deg r\left(x\right)<\deg \left(x-b\right)=1\Longrightarrow r\left(x\right)$为非零常数,
令上式$x=b\Longrightarrow r\left(x\right)=f\left(b\right)$.
\end{Proof}
\end{formal}
\begin{formal}
\begin{definition}[重根]\label{def:重根}
    设$f\left(x\right)\in\mathbb{K}\left[x\right],b\in\mathbb{K}$,若存在
    $k\in\mathbb{Z}^+$\,\textup{s.t.}$\left(x-b\right)^k\mid f\left(x\right)$但$\left(x-b\right)^{k+1}\nmid f\left(x\right).$称$b$是$f\left(x\right)$的$k$重根.特别地,$k=1$时称为单根.  
\end{definition}
\end{formal}
\begin{formal}
\begin{definition}[等价定义]\label{def:等价定义}
    上面$k$重根等价于
    \[
    f\left(b\right)=f'\left(b\right)=\cdots=f^{\left(k-1\right)}\left(b\right)=0,f^{\left(k\right)}\left(b\right)\neq 0
    \]
\end{definition}
\end{formal}
\begin{red}
\begin{remark}
    $k$重根视作$k$个根.
\end{remark}
\end{red}
\begin{green}
\begin{lemma}[不可约多项式的根]\label{lem:不可约多项式的根}
    设$\mathbb{K}$上的不可约多项式$f\left(x\right)$,若$\deg f\left(x\right)\geqslant 2$,则$f\left(x\right)$在$\mathbb{K}$上没有根.
\end{lemma}
\begin{Proof}
考虑反证法,设$b\in\mathbb{K}$是$f\left(x\right)$的根,从而$\left(x-b\right)\mid f\left(x\right)$即
\[
f\left(x\right)=\left(x-b\right)q\left(x\right)
\]
于是$\deg q\left(x\right)\geqslant 1.$这与不可约矛盾.证毕.
\end{Proof}
\end{green}
\begin{formal}
\begin{theorem}[根的个数]\label{thm:根的个数}
    设数域$\mathbb{K}$上的$n$次多
    项式$f\left(x\right)$,
    则$f\left(x\right)$在数域$
    \mathbb{K}$上至多有$n$个根.
\end{theorem}
\begin{Proof}
作标准因式分解
\[
f\left(x\right)=cp_1\left(x\right)^{e_1}p_2\left(x\right)^{e_2}\cdots p_m\left(x\right)^{e_m}
\]
将一次项单独拿出
\[
f\left(x\right)=c\left(x-b_1\right)^{n_1}\left(x-b_2\right)^{n_2}\cdots\left(x-b_r\right)^{n_r}p_1\left(x\right)^{e_1}p_2\left(x\right)^{e_2}\cdots p_s\left(x\right)^{e_s}
\]
其中$c\neq 0,b_1,b_2,\cdots,b_r$互不相同,$p_1\left(x\right),p_2\left(x\right),\cdots,p_s\left(x\right)$是互异的首一不可约多项式且$\deg p_i\left(x\right)\geqslant 1\left(\forall 1\leqslant i\leqslant 
s\right)$.于是$f\left(x\right)$的根为$b_1,b_2,\cdots,b_r$,其分别为$n_1,n_2,\cdots,n_r$重.两边同时取次数即得
\[
n=n_1+n_2+\cdots+n_r+e_1\deg p_1\left(x\right)+e_2\deg p_2\left(x\right)+\cdots+e_s\deg p_s\left(x\right)
\]于是$f\left(x\right)$在数域$\mathbb{K}$上的根一共有$n_1+n_2+\cdots+n_r\leqslant n$个.
\end{Proof}
\end{formal}
\begin{green}
\begin{lemma}[多项式相同]\label{lem:多项式相同}
    设多项式$f\left(x\right),g\left(x\right)\in\mathbb{K}\left[x\right],\deg f\left(x\right)\leqslant n,\deg g\left(x\right)\leqslant n$,若$\mathbb{K}$中存在$n+1$个数$b_1,b_2,\cdots,b_{n+1}$\,\textup{s.t.}
    \[
    f\left(b_i\right)=g\left(b_i\right)\left(\forall 1\leqslant i\leqslant n+1\right)
    \]
    则$f\left(x\right)=g\left(x\right)$（作为多项式而言）.
\end{lemma}
\begin{Proof}
设$h\left(x\right)=f\left(x\right)-g\left(x\right)$,由\cref{thm:根的个数}显然.
\end{Proof}
\end{green}
至此,我们得到了作为多项式相等则作为多项式函
数相等的结论.下面证明作为多项式函数相等则作
为多项式相等.
\begin{brown}
\begin{Proof}
设$\deg f\left(x\right)\leqslant n,\deg g\left(x\right)\leqslant n.$因为$\mathbb{K}\supseteq \mathbb{Q}$,故能取出$n+1$个不同的数
$b_1,b_2,\cdots,b_{n+1}\in\mathbb{K}$
\textup{s.t.}$f\left(b_i\right)=g\left(b_i\right)\left(\forall 1\leqslant i\leqslant n+1
\right)$,则二者作为多项式恒等.
\end{Proof}
\end{brown}
\subsection{域}
\begin{red}
\begin{remark}
    上述证明的关键在
    于“数域”$\mathbb{K}$是无穷集.但
    对于一般的域而言,却不见得成立
    .
    
    在一般的域上,一般来说,作为多项
    式相等和作为多项式函数相等,是不一样的.
\end{remark}
\end{red}
\begin{red}
\begin{remark}
    再次强调,数域只是域的一种.下面将描述一般
    的域上的该问题.
\end{remark}
\end{red}
\begin{brown}
\begin{example}
    考虑伽罗瓦域$\mathbb{F}_2=\left\{\overline{0},\overline{1}\right\}$,这里的$\overline{0}$可
    以理解为全体偶数,$\overline{1}$理解为全体奇
    数.定义加法
    \begin{align*}
        \overline{0}+\overline{0}=\overline{0},\overline{1}+\overline{1}=\overline{0},\overline{1}+\overline{0}=\overline{1}
    \end{align*}
    和乘法
    \[
    \overline{0}\cdot \overline{0}=\overline{0},
    \overline{0}\cdot \overline{1}=\overline{0},
    \overline{1}\cdot\overline{1}=\overline{1}
    \]
    容易验证这个域对加减乘除封闭.同时这是一个有限域,那么上面的定理就不成立了.

    比如,考虑$f\left(x\right)=x^2,g\left(x\right)=x$,二者作为多项式不相等,
    但是作为$\mathbb{F}_2$上的
    多项式函数是相等的.
\end{example}
\end{brown}
\begin{red}
\begin{remark}
    高等代数仅涉及数域,不会涉及更加一般的域.
\end{remark}
\end{red}
\newpage
\section{复系数多项式}
\subsection{二元连续函数}
我们需要二元连续函数的一些相关知识作为铺垫.
\begin{formal}
\begin{definition}[连续]\label{def:连续}
    设$f\left(x,y\right)$是$\mathbb{R}^2$上的函数,
    则$f\left(x,y\right)$在$\left(x_0,y_0\right)$
    连续定义为$\forall \epsilon >0,\exists \delta >0,\forall \left(x,y\right)\in O_{\left(x_0,y_0\right),\delta}$
    \[
    \left|f\left(x,y\right)-f\left(x_0,y_0\right)\right|<\epsilon
    \]
\end{definition}
\end{formal}
\begin{formal}
\begin{theorem}[二元多项式连续性]\label{thm:二元多项式连续性}
    二元多项式都是二元连续函数.
\end{theorem}
\end{formal}
\begin{formal}
\begin{proposition}[二元连续函数的性质]\label{prop:二元连续函数的性质}
    任一二元连续函数$f\left(x,y\right)$一定能在闭圆盘$D=\left\{
        \left(x,y\right)\big|\left(x-
        x_0\right)^2+\left(y-y_0\right)^2\leqslant r^2
    \right\}$上取到最大值和最小值.
\end{proposition}
\end{formal}
在承认该命题的前提下,我们来证明代数学基本定理.
\subsection{代数学基本定理}
\begin{formal}
\begin{theorem}[代数学基本定理]\label{thm:代数学基本定理}
    任何次数大于等于$1$的复系数多项式至少有一个复根.或者说,复数域上的不可约多项式
    只能是一次的.
\end{theorem}
\begin{Proof}
设$f\left(z\right)=a_nz^n+a_{n-1}z^{n-1}+\cdots+a_1z+a_0\in\mathbb{C}\left[z\right],a_n\neq 0,n\geqslant 1.
$

第一步,证明：$\exists z_0\in\mathbb{C}$\,\textup{s.t.}$\forall z\in\mathbb{C},
\left|f\left(z\right)\right|\geqslant \left|f\left(z_0\right)\right|$.

设$z=x+y\mathrm{i},x,y\in\mathbb{R}$,则$f\left(z\right)=f\left(x+y\mathrm{i}\right)=
P\left(x,y\right)+Q\left(x,y\right)\mathrm{i}$,其中$P\left(x,y\right),Q\left(x,y\right)
$是二元多项式.那么\[
\left|f\left(z\right)\right|
=\sqrt{P\left(x,y\right)^2+
Q\left(x,y\right)^2}
\]
故$\left|f\left(z\right)\right|$是关于$x,y$
的二元连续函数.

考虑估计
\[
\left|z_1\right|-\left|z_2\right|\leqslant \left|z_1+z_2\right|\leqslant \left|z_1\right|+\left|z_2\right|,\forall z_1,z_2\in\mathbb{C}
\]
那么
\begin{align*}
    \left|f\left(z\right)\right|&=\left|a_nz^n+a_{n-1}z^{n-1}+\cdots+a_1z+a_0\right|\\
&\geqslant \left|a_n\right|\left|z\right|^n-\left|a_{n-1}z^{n-1}+\cdots+a_1z+a_0\right|\\
&\geqslant\cdots\\
&\geqslant \left|z\right|^n\left(
    \left|a_n\right|-
    \left(
        \frac{\left|a_{n-1}\right|}{\left|z\right|}+
        \cdots+
        \frac{\left|a_1\right|}{\left|z^{n-1}\right|}+
        \frac{\left|a_0\right|}{\left|z^n\right|}
    \right)
\right)
\end{align*}
$\left|z\right|\to \infty$时,上式$\displaystyle
\geqslant \frac{1}{2}\left|a_n\right|\left|z\right|^n\to
 \infty.$

 任取$z_1\in\mathbb{C},\exists R\in\mathbb{R},\textup{s.t.}\left|f\left(z\right)\right|\geqslant \left|f\left(z_1\right)\right|\left(
    \left|z\right|>R\right)$,由二元连续函数的性质知,
    $\left|f\left(z\right)\right|$在闭圆盘
    $D=\left\{z\mid \left|z\right|\leqslant R\right\}$
    取得最小值.不妨设为$z_0$,即
    \[
    \left|f\left(z\right)\right|\geqslant \left|f\left(z_0\right)\right|,\forall z\in 
    D
    \]
    于是$\left|f\left(z\right)\right|\left(z\in\mathbb{C}
    \right)$在$z_0$取
    到最小值.

    第二步,证明$f\left(z_0\right)=0.$考虑反证法,设$f\left(z_0\right)\neq 0.$
    作一微小摄动$z=z_0+h$,则
    \[
    f\left(z_0+h\right)=b_nh^n+b_{n-1}h^{n-1}+\cdots+b_1h+b_0
    \]
    其中,$b_n=a_n\neq 0,b_0=f\left(z_0\right)\neq 0.$则
    \begin{align*}
         \frac{f\left(z_0+h\right)}{f\left(z_0\right)}&=1+\frac{b_1}{f\left(z_0\right)}h+
    \cdots+\frac{b_{n-1}}{f\left(z_0\right)}h^{n-1}+\frac{b_n}{f\left(z_0\right)}h^n\\
    &=1+c_1h+\cdots+c_{n-1}h^{n-1}+c_nh^n
    \end{align*}
    设存在$1\leqslant k\leqslant n\,\textup{s.t.}\,c_1=\cdots=c_{k-1}=0,c_k\neq 0.
    $于是
    \begin{align*}
        \frac{f\left(z_0+h\right)}{f\left(z_0\right)}=1+c_kh^k+\cdots+c_nh^n
    \end{align*}
    令$\displaystyle
    h=e\cdot d,d=\sqrt[k]{-\frac{1}{c_k}},0<e\ll 1.$于是
    \begin{align*}
     \left|    
        \frac{f\left(z_0+h\right)}{f\left(z_0\right)}\right|&=
    \left|1-e^k+c_{k+1}e^{k+1}d^{k+1}+\cdots+c_ne^nd^n
    \right|
    \\
    &\leqslant 
    \left|1-e^k\right|+\left|c_{k+1}e^{k+1}d^{k+1}+\cdots+c_ne^nd^n
    \right|\\
    &=1-e^k+e^{k+1}\left(\left|c_{k+1}d^{k+1}\right|+\cdots+\left|c_nd^n\right|\right)
    \end{align*}
    当$e$充分小时,$\displaystyle
    LHS\leqslant 1-e^k+\frac{1}{2}e^k<1$.这与$\left|f\left(z_0\right)\right|$最小矛盾.

    证毕.
   \end{Proof}
\end{formal}
\begin{red}
\begin{remark}
    代数学基本定理有很多等价的命题,证明方法也很多.
\end{remark}
\end{red}
\begin{green}
\begin{corollary}[代数学基本定理的等价命题]\label{cor:代数学基本定理的等价命题}
    设$n\geqslant 1$,则下列结论等价
    \begin{enumerate}[label={\textup{(\arabic*)}}]
        \item 任一$n$次复系数多项式至少有一个复根
        \item 复数域上的不可约多项式都是一次多项式
        \item 任一复系数多项式都是一次多项式的乘积
        \item 任一$n$次复系数多项式恰好有$n$个复根（计重数）
    \end{enumerate}
\end{corollary}
\begin{Proof}
$(1)\Longrightarrow (2)$设$p\left(x\right)$在$\mathbb{C}$不可约,
设$p\left(b\right)=0$则由前文知$\left(x-b\right)|p\left(x\right)\Longrightarrow p\left(x\right)=\left(x-b\right)q\left(x\right)\Longrightarrow q\left(x\right)=c\neq 0.$

$(2)\Longrightarrow (3)$由因式分解定理立得.

$(3)\Longrightarrow (4)$$f\left(x\right)=c\left(x-b_1\right)\left(x-b_2\right)\cdots\left(x-b_n\right)$.则$(4)\Longrightarrow (1)$显然.
\end{Proof}
\end{green}
\subsection{Vieta定理}
\begin{formal}
\begin{theorem}[Vieta定理]\label{thm:Vieta定理}
    设$f\left(x\right)=a_0x^n+a_1x^{n-1}+\cdots+a_{n-1}x+a_n\in\mathbb{K}\left[x
    \right]$且其复根为$x_1,x_2,\cdots,x_n$,则
    \[
    \sum_{i=1}^{n}x_i=-\frac{a_1}{a_0},
    \]
    \[
    \sum_{1\leqslant i<j\leqslant n}x_ix_j=\frac{a_2}{a_0},
    \]
    \[
    \sum_{1\leqslant i<j<k\leqslant n}x_ix_jx_k=-\frac{a_3}{a_0},
    \]
    \[
    \cdots\cdots\cdots\cdots
    \]
    \[
    \sum_{1\leqslant i_1<i_2<\cdots<i_r\leqslant n}x_{i_1}x_{i_2}\cdots x_{i_r}=\left(-1\right)^r\frac{a_r}{a_0},
    \]
    \[
    \cdots\cdots\cdots\cdots
    \]
    \[
    \prod_{i=1}^{n}x_i=\left(-1\right)^n\frac{a_n}{a_0}
    \]
\end{theorem}
\begin{Proof}
$f\left(x\right)=c\left(x-x_1\right)\left(x-x_2\right)\cdots\left(x-x_n\right)$,展开即得.
\end{Proof}
\end{formal}
\subsection{一元方程的求解}
一元$n$次方程的求解在历史上是持久的问题,下面都是在复数域$\mathbb{C}$讨论.
\subsubsection{一元二次方程}
一元二次方程$ax^2+bx+c=0$的求根公式为
\[
x_{1,2}=\frac{
    -b\pm \sqrt{b^2-4ac}
}{2a}
\]
\subsubsection{一元三次方程}
设方程
\[
x^3+ax^2+bx+c=0
\]
作$\displaystyle
x=y-\frac{1}{3}a$得
\[
y^3+py+q=0
\]

于是问题转为对形如
$f\left(x\right)=x^3+px+q=0$的方程的解的讨论.

\begin{enumerate}[label={\textup{(\arabic*)}}]
    \item $q=0$时$x_1=0,x_2=\sqrt{-p},x_3=-\sqrt{-p}.p=0$时$x_1=\sqrt[3]{-q},x_2=\sqrt[3]{-q}\omega,x_3=\sqrt[3]{-q}\omega^2.$其中,$\omega$为三次单位根
    \[
    \omega ^3=1,\omega = -\frac{1}{2}+\frac{\sqrt{3}}{2}\mathrm{i}
    \]
    因此解的形式也可以是
    $x_i=\sqrt[3]{-q}\omega^i\left(i=1,2,3
    \right)$.
    \item 考虑$p\neq 0,q\neq 0$.引$x=u+v$,则
    \[
    x^3-3uvx-\left(u^3+v^3\right)=0
    \]
    于是\[
    \begin{cases*}
    uv=-\frac{1}{3}p\\
    u^3+v^3=-q
    \end{cases*}
    \]
    则\[
    u^3,v^3=-\frac{q}{2}\pm \sqrt{\frac{q^2}{4}+\frac{p^3}{27}}
    \]
    令\[
    \Delta = \frac{q^2}{4}+\frac{p^3}{27}
    \]
    则
    \[
    \begin{cases*}
    x_1=\sqrt[3]{-\cfrac{q}{2}+\sqrt{\Delta}}+\sqrt[3]{-\cfrac{q}{2}-\sqrt{\Delta}}\\
    x_2=\omega \sqrt[3]{-\cfrac{q}{2}+\sqrt{\Delta}}+\omega^2\sqrt[3]{-\cfrac{q}{2}-\sqrt{\Delta}}\\
    x_3=\omega^2 \sqrt[3]{-\cfrac{q}{2}+\sqrt{\Delta}}+\omega\sqrt[3]{-\cfrac{q}{2}-\sqrt{\Delta}}
    \end{cases*}
    \]
    此式称为\textup{Cardano}（卡尔达诺公式）
\end{enumerate}
\subsubsection{一元四次方程——Ferrari解法}
考虑\[
x^4+ax^3+bx^2+cx+d=0
\]
作$\displaystyle
x=y-\frac{1}{4}a$转为求解方程
\[
x^4+ax^2+bx+c=0
\]
引未知数$u$并加减$\displaystyle
ux^2+\frac{u^2}{4}$项得
\[
\left(x^2+\frac{u}{2}\right)^2=\left[
    \left(u-a\right)x^2-bx+\frac{u^2}{4}-c
\right]
\]
右边是完全平方的条件为
\[
b^2-4\left(u-a\right)\left(\frac{u^2}{4}-c\right)=0
\]
称为预解式.解出$u$后（考虑到对称性,此三次方程的根取其一即可
）得\[
\left(x^2+\frac{u}{2}\right)^2=
\left(
    \sqrt{u-a}x-\frac{b}{2\sqrt{u-a}}
\right)^2
\]
得到两个二次方程,求解即可.

这一求解方法被称为\textup{Ferrari}（费拉里
）解法.
\subsubsection{更高次方程}
数学家们早期发现高于四次的方程没有一般形式的根式解,但无法证明.这一结论在
19世纪30年代被天才数学家\textup{Galois}（伽罗瓦）证明.即五次及以上的方程没有求根公式,但并非不可解.

对于具体的一个方程,
Galois也给出了一套判定,即一个多项式的根可以以根式形式表示,等价于
这个多项式的Galois群是可解群.
\newpage
\section{实系数多项式与有理系数多项式}
\begin{red}
    \begin{remark}
因为\[
\mathbb{Q}\subseteq \mathbb{R}\subseteq\mathbb{C}
\]
我们来分别考虑三个数域上的不可约多项式.
    \end{remark}
\end{red}
\subsection{实数域上的不可约多项式}
我们已经证明复数域$\mathbb{C}$上的不可约多项式只能是一次多项式,那么实数域上的不可约多项式又是什么？
\begin{formal}
\begin{theorem}[根的成对]\label{thm:根的成对}
    设$f\left(x\right)=a_nx^n+a_{n-1}x^{n-1}+\cdots+a_1x+a_0\in\mathbb{R}\left[x\right]$,若$a+b\mathrm{i}\left(b\neq 0\right)$是一个复根,
    那么$a-b\mathrm{i}$也是一个根.
\end{theorem}
\begin{Proof}
利用$\forall z_1,z_2\in\mathbb{C},\overline{z_1+z_2}=\overline{z_1}+\overline{z_2},\overline{z_1\cdot z_2}=\overline{z_1}\cdot \overline{z_2}$即可.
\end{Proof}
\end{formal}
\begin{formal}
\begin{theorem}[实数域上的不可约多项式]\label{thm:实数域上的不可约多项式}
    实数域上的不可约多项式只能是一次多项式或者如下无实数根的二次多项式：
    \[
    ax^2+bx+c=0,\Delta = b^2-4ac<0
    \]
\end{theorem}
\begin{Proof}
任取$\mathbb{R}$上的不可约多项式,
其一,$\deg f\left(x\right)=1$显然成立.

其二,
$\deg f\left(x\right)=2$时设$f\left(x\right)=ax^2+bx+c\left(a\neq 0\right)$,
则$f\left(x\right)$可约当且仅当$f\left(x\right)=a\left(x-x_1\right)\left(x-x_2\right)
\left(x_1,x_2\in\mathbb{R}\right)$.于是$f\left(x\right)$不可约当且仅当其无实根即判别式
小于零.

其三,考虑$\deg f\left(x\right)\geqslant 3$.
一方面,如果$f\left(x\right)$有实根$x_0$,
那么$f\left(x\right)=\left(x-x_0\right)q\left(x\right)\Longrightarrow $可约.
另一方面,如果有虚根,一定成对,那么\[\left(
    x-\left(a+b\mathrm{i}\right)
\right)\big|_{\mathbb{C}}f\left(x\right),
\left(
    x-\left(a-b\mathrm{i}\right)
\right)\big|_{\mathbb{C}}f\left(x\right)
\]
因为$x-\left(a+b\mathrm{i}\right)$和$x-\left(a-
b\mathrm{i}\right)$均是不可约的且互异,于是互素,于是
\[
\left(\left(x-a\right)^2+b^2\right)\big|_{\mathbb{C}}f\left(x
\right)\Longrightarrow
\left(\left(x-a\right)^2+b^2\right)\big|_{\mathbb{R}}f\left(x\right)
\]
（因为整除关系在基域扩张下不
改变\textsuperscript{\cite{谢启鸿2015}}）
即$f\left(x\right)=g\left(x\right)\left(\left(x-a\right)^2+b^2\right),
g\left(x\right)\in\mathbb{R}\left[x\right]\Longrightarrow \deg g\left(x\right)\geqslant 1\Longrightarrow $可约.

另证\[
\left(\left(x-a\right)^2+b^2\right)\big|_{\mathbb{R}}f\left(x
\right)
\]
考虑$\mathbb{R}$上的带余除法
\[
f\left(x\right)=\left(\left(x-a\right)^2+b^2\right)q\left(x\right)+r\left(x\right)
\]
$q\left(x\right),r\left(x\right)\in\mathbb{R}\left[x\right],\deg r\left(x\right)<\deg \left(\left(x-a\right)^2+b^2\right)=2$
.代入$x=a+b\mathrm{i}$得$0=r\left(a+b\mathrm{i}\right)$,结合次数即得
$r\left(x\right)=0.$
\end{Proof}
\end{formal}
\begin{green}
\begin{corollary}[因式分解定理]\label{cor:因式分解定理}
    实数域上的多项式$f\left(x\right)$必定可以分解为
    有限个一次因式及不可约二次因式（判别式小于零）的乘积.
\end{corollary}
\end{green}
\subsection{有理数域上的不可约多项式}
\subsubsection{整系数多项式与有理系数多项式}
\begin{formal}
\begin{theorem}[整系数多项式存在有理根的必要条件]\label{thm:整系数多项式存在有理根的必要条件}
    设$n$次整系数多项式
    \[
    f\left(x\right)=a_nx^n+a_{n-1}x^{n-1}+\cdots+a_1x+a_0
    \]
    其中$a_n\neq 0,p,q$为互素整数,则有理数$\displaystyle
    \frac{p}{q}$是$f\left(x\right)$的根的必要条件是$q\mid a_n,p\mid a_0.$
\end{theorem}
\begin{Proof}
令$\displaystyle
x=\frac{p}{q}$即可.
\end{Proof}
\end{formal}
\begin{red}
\begin{remark}
    因为任一有理系数方程均可化为同解的整系数方程,因此这也是
有理系数方程是否存在有理根的必要条件.
\end{remark}
\end{red}
\begin{brown}
\begin{example}
$    f\left(x\right)=x^5-12x^3+36x+12=0$无有理根.
\end{example}
\end{brown}
\begin{red}
\begin{remark}
    设$f\left(x\right)\in\mathbb{Q}\left[x\right],\deg f\left(x\right)\geqslant 2$,则
    \[
    f\left(x\right)\text{有有理根}\Longrightarrow
    f\left(x\right)\text{在}\mathbb{Q}\text{可约}
    \]
    反推不成立.
\end{remark}
\end{red}
\begin{red}
\begin{remark}
    有理系数多项式在有理数域上的可约性与不可约性的研究完全等价于其化为的整系
    数多项式的可约性与不可约性的研究.
\end{remark}
\end{red}
\subsubsection{本原化}
\begin{formal}
\begin{definition}[本原多项式]\label{def:本原多项式}
    设整系数多项式$f\left(x\right)=a_nx^n+a_{n-1}x^{n-1}+\cdots+a_1x+a_0$
    ,若$\mathrm{g.c.d.}\left(a_n,a_{n-1},
    \cdots,a_1,a_0\right)=1$,则称$f\left(x\right)$
    为本原多项式.
\end{definition}
\end{formal}
\begin{brown}
\begin{example}
    首一多项式均为本原多项式.
\end{example}
\end{brown}
\subsubsection{可约的延拓}
\begin{green}
\begin{lemma}[Gauss引理]\label{lem:Gauss引理}
    两个本原多项式之积仍是本原多项式.
\end{lemma}
\begin{Proof}
设本原多项式
\begin{align*}
    &f\left(x\right)=a_nx^n+a_{n-1}x^{n-1}+\cdots+a_1x+a_0;\\
&g\left(x\right)=b_mx^m+b_{m-1}x^{m-1}+\cdots+b_1x+b_0
\end{align*}
积\[
f\left(x\right)g\left(x\right)=c_{n+m}x^{n+m}+
c_{n+m-1}x^{n+m-1}+\cdots+c_1x+c_0
\]
其中$\displaystyle
c_k=\sum_{i+j=k}a_ib_j.$考虑反证法,设$f\left(x\right)g\left(x\right)$不是本原多项式,
则存在一个素数$p$\,\textup{s.t.}$p\mid c_k\left(\forall 0\leqslant 
k\leqslant n+m \right)$.因为$f\left(x\right)$是本原多项式,于是$\exists 0\leqslant i\leqslant n\,\textup{s.t.}
p\mid a_0,p\mid a_1,\cdots,p\mid a_{i-1},p\nmid a_i.$同理,
$\exists 0\leqslant j\leqslant m\,\textup{s.t.}p\mid b_0,p\mid b_1,\cdots
,p\mid b_{j-1},p\nmid b_j.$因为
\[
c_{i+j}=\cdots+a_{i-2}b_{j+2}+a_{i-1}b_{j+1}+a_ib_j+a_{i+1}b_{j-1}+a_{i+2}b_{j-2}+\cdots
\]
那么一定有$p\mid a_ib_j$矛盾.
\end{Proof}
\end{green}
\begin{formal}
\begin{definition}[整数环上的可约与不可约]\label{def:整数环上的可约与不可约}
    若一个整系数多项式可以分解为两个次数较低的整系数多项式之积,则称为在整数环上可约.
\end{definition}
\end{formal}
\begin{formal}
\begin{theorem}[可约范围的缩小]\label{thm:可约范围的缩小}
    设$f\left(x\right)$为整系数多项式,则$f\left(x\right)$在$\mathbb{Q}$可约等价于$f\left(x\right)$在$\mathbb{Z}$可约.即$f\left(x\right)=g\left(x\right)
    h\left(x\right),g\left(x\right),h\left(x\right)\in\mathbb{Z}\left[x\right]$,其中$\deg\,g\left(x\right)\geqslant 1,\deg\,h\left(x\right)\geqslant 1.$
\end{theorem}
\begin{Proof}
先证充分性,即$f\left(x\right)=g\left(x\right)h\left(x\right),g\left(x\right),h\left(x\right)\in\mathbb{Z}\left[x\right],\deg \,g\left(x\right)\geqslant 1,
\deg\,h\left(x\right)\geqslant 1$,将$g\left(x\right),h\left(x\right)$视为
$\mathbb{Q}$上的多项式,则$f\left(x\right)$在$\mathbb{Q}$可约.

然后证必要性,设$f\left(x\right)$在$\mathbb{Q}$可约,即\[
f\left(x\right)=g\left(x\right)h\left(x\right)
\]
其中$g\left(x\right),h\left(x\right)\in\mathbb{Q}\left[x\right],\deg\,g\left(x\right)\geqslant 1,\deg\,h\left(x\right)\geqslant1.$设$g\left(x\right)=ag_1\left(x\right),a\in\mathbb{Q},g_1\left(x\right)$是
本原的整系数多项式,$h\left(x\right)=bh_1\left(x\right),b\in\mathbb{Q},h_1\left(x\right)$是本原的整系数多项式.于是
\[
f\left(x\right)=abg_1\left(x\right)h_1\left(x\right)
\]
断言$ab\in\mathbb{Z}$.考虑反证法,设$\displaystyle
ab=\frac{p}{q},\left(p,q\right)=1$且$q>1.$于是$\displaystyle
\frac{p}{q}g_1\left(x\right)h_1\left(x\right)\notin \mathbb{Z}\left[x\right]$矛盾.于是$ab\in\mathbb{Z}$则
\[
f\left(x\right)=\left(abg_1\left(x\right)\right)h_1\left(x\right)
\]
证毕.
\end{Proof}
\end{formal}
\begin{formal}
\begin{theorem}[等价表述]\label{thm:等价表述}
    \cref{thm:可约范围的缩小}的等价表述为：
    设$f\left(x\right)\in\mathbb{Z}\left[x\right]$,则$f\left(x\right)$在$\mathbb{Z}$不可约等价于$f\left(x\right)$在$\mathbb{Q}$不可约.
\end{theorem}
\end{formal}
\subsubsection{Eisenstein判别法}
\begin{formal}
\begin{theorem}[Eisenstein判别法]\label{thm:Eisenstein判别法}
    设整系数多项式$f\left(x\right)=a_nx^n+a_{n-1}x^{n-1}+\cdots+a_1x+a_0,a_n\neq 0,n\geqslant 1,p$为素数.若$p\mid a_i\left(\forall 0\leqslant i\leqslant n-1\right)$但$p\nmid a_n$且$p^2
    \nmid a_0$则$f\left(x\right)$
    在$\mathbb{Q}$不可约.
\end{theorem}
\begin{proof}
    只需要证明在$\mathbb{Z}$不可约即可.考虑反证法,
    设$f\left(x\right)=g\left(x\right)h\left(x\right)$,其中
    \newpage
    \begin{align*}
        g\left(x\right)=b_mx^m+b_{m-1}x^{m-1}+\cdots+b_1x+b_0\\
        h\left(x\right)=c_tx^{t}+c_{t-1}x^{t-1}+\cdots+c_1x+c_0
    \end{align*}
    是整系数多项式且$m+t=n,
    m\geqslant1,t\geqslant 1,a_n=b_mc_t,a_0=b_0c_0.$因为
    $p\nmid a_n\Longrightarrow p\nmid b_m$且
    $p\nmid c_t.$
    又$p\mid a_0=b_0c_0
    \Longrightarrow p\mid b_0$或$p\mid c_0$且不能同时成立$\left(p^2\nmid a_0
    \right)$.不妨设$p\mid b_0,p\nmid c_0$则$\exists 1\leqslant j\leqslant m,p\mid b_0,p\mid b_1,\cdots,p\mid b_{j-1},p\nmid 
    b_j$,因为\[
    a_j=b_0c_j+b_1c_{j-1}+\cdots+b_{j-1}c_1+b_jc_0
    \]
    于是$p\nmid a_j$,矛盾.
\end{proof}
\end{formal}
\begin{red}
\begin{remark}
    \textup{Eisenstein}
    判别法是一个非常强的充分条件,但是正如我们前面所说,
    目前没有比较好的充分必要条件.
\end{remark}
\end{red}
\begin{brown}
\begin{example}
    证明：$\forall n\geqslant 1,x^n-2$在$\mathbb{Q}$不可约.
\end{example}
\begin{Proof}
\cref{thm:Eisenstein判别法}取$p=2$即可.
\end{Proof}
\end{brown}
\begin{red}
\begin{remark}
    在$\mathbb{Q}$上,不可约多项式的次数没有上界.
\end{remark}
\end{red}
\begin{brown}
\begin{example}
    设素数$p$,证明：
    \[
    f\left(x\right)=x^{p-1}+x^{p-2}+\cdots+x+1
    \]
    在$\mathbb{Q}$不可约.
\end{example}
\begin{Proof}
考虑$x=y+1$则\begin{align*}
    f\left(x\right)=\frac{x^p-1}{x-1}=\frac{\left(y+1\right)^p-1}{y}=
    y^{p-1}+\mathrm{C}_p^1y^{p-2}+\mathrm{C}_p^2y^{p-3}+\cdots+\mathrm{C}_p^{p-1}
\end{align*}
注意到$p\mid \mathrm{C}_p^i\left(1\leqslant i\leqslant p-1\right),p\nmid 1,p\nmid \mathrm{C}_p^{p-1}=p$
,由\cref{thm:Eisenstein判别法}证毕.
\end{Proof}
\end{brown}
\begin{brown}
\begin{example}
    证明
    \[
    f\left(x\right)=1+x+\frac{x^2}{2!}+\cdots+\frac{x^n}{n!}
    \]
    在$n$为素数时在$\mathbb{Q}$不可约.
\end{example}
\begin{Proof}
考虑$n!f\left(x\right),n=p$由\cref{thm:Eisenstein判别法}即可.
\end{Proof}
\end{brown}
\subsubsection{Osada判别法}
\begin{formal}
\begin{theorem}[Osada判别法]\label{thm:Osada判别法}
    设首一整系数多项式
    \[
    f\left(x\right)=x^n+a_{n-1}x^{n-1}+\cdots+a_1x+a_0
    \]
    其中$\displaystyle
    \left|a_0\right|>1+\sum_{i=1}^{n-1}\left|a_i\right|,\left|a_0\right|$
    为素数.则$f\left(x\right)$在$\mathbb{Q}$不可约.
\end{theorem}
\end{formal}
\section{多元多项式}
\subsection{定义}
\begin{formal}
\begin{definition}[单项式]\label{def:单项式}
    设$\mathbb{K}$为数域,$x_1,x_2,\cdots,x_n$是未定元,则
    \[
    ax_1^{i_1}x_2^{i_2}\cdots x_n^{i_n}
    \]
    （其中,$a\in\mathbb{K}$称为系数,$i_1,i_2,\cdots,i_n\in\mathbb{N}$）称为
    一个单项式.

    单项式的次数定义为
    \[
    \deg\,\left(ax_1^{i_1}x_2^{i_2}\cdots x_n^{i_n} \right):=
    \begin{cases*}
    -\infty &,$a=0$\\
    i_1+i_2+\cdots+i_n&,$a\neq 0$
    \end{cases*}
    \]
\end{definition}
\end{formal}
\begin{formal}
\begin{definition}[同类项]\label{def:同类项}
    若 $ax_1^{i_1}x_2^{i_2}\cdots x_n^{i_n}\left(a\neq 0\right)$ 与 $bx_1^{j_1}x_2^{j_2}\cdots x_n^{j_n}\left(b\neq 0\right)$ 为同类项,则 $i_1=j_1,i_2=j_2,\cdots,i_n=j_n$.
\end{definition}
\end{formal}
\begin{formal}
\begin{definition}[合并同类项]\label{def:合并同类项}
    定义两个同类项之间的加法就是系数的加法,即若 $ax_1^{i_1}x_2^{i_2}\cdots x_n^{i_n}$ 与 $bx_1^{i_1}x_2^{i_2}\cdots x_n^{i_n}$ 为同类项,则
    \[ ax_1^{i_1}x_2^{i_2}\cdots x_n^{i_n}+bx_1^{i_1}x_2^{i_2}\cdots x_n^{i_n}=\left(a+b\right)x_1^{i_1}x_2^{i_2}\cdots x_n^{i_n} \]
\end{definition}
\end{formal}
\begin{formal}
\begin{definition}[单项式的运算定义]\label{def:单项式的运算定义}
    同类项的加法即是合并同类项,非同类项的加法即是形式和；$\forall k \in\mathbb{K}$,数乘定义为
    \[ kax_1^{i_1}x_2^{i_2}\cdots x_n^{i_n}=\left(ka\right)x_1^{i_1}x_2^{i_2}\cdots x_n^{i_n} \]
    乘法定义为
    \[ ax_1^{i_1}x_2^{i_2}\cdots x_n^{i_n}\cdot bx_1^{j_1}x_2^{j_2}\cdots x_n^{j_n}=\left(ab\right)x_1^{i_1+j_1}x_2^{i_2+j_2}\cdots x_n^{i_n+j_n} \]
\end{definition}
\end{formal}
\begin{formal}
\begin{definition}[多元多项式]\label{def:多元多项式}
    有限个单项式的形式和称为多元多项式即
    \[
    f\left(x_1,x_2,\cdots,x_n\right)
    =
    \sum_{\left(i_1,i_2,\cdots,i_n\right)\in \mathbb{N}^n}a_{i_1i_2\cdots i_n}x_1^{i_1}x_2^{i_2}\cdots x_n^{i_n}\left(\textup{有限形式和}\right)
    \]
    其次数定义为
    \[
    \deg\,f\left(x\right):=
    \begin{cases*}
    \max\left\{
    i_1+i_2+\cdots+i_n\mid a_{i_1i_2\cdots i_n}\neq 0
    \right\}&,$f\left(x\right)\neq 0$\\
    -\infty&,$f\left(x\right)=0$
    \end{cases*}
    \]
\end{definition}
\end{formal}
\begin{formal}
\begin{definition}[多元多项式的相等的刻画]\label{def:多元多项式的相等的刻画}
    两个多元多项式 $\displaystyle
    f\left(x_1,x_2,\cdots,x_n\right)=\sum
    _{\left(i_1,i_2,\cdots,i_n\right)} 
    a_{i_1i_2\cdots i_n}x_1^{i_1}x_2^{i_2}\cdots x_n^{i_n}$ 与 $\displaystyle
    g\left(x_1,x_2,\cdots,x_n\right)=
    \sum_{\left(i_1,i_2,\cdots,i_n\right)}b_{i_1i_2\cdots i_n}x_1^{i_1}x_2^{i_2}\cdots x_n^{i_n}$ 相等当且仅当$a_{i_1i_2\cdots i_n}=b_{i_1i_2\cdots i_n},\forall \left(i_1,i_2,\cdots,i_n\right)\in\mathbb{N}^n$.
\end{definition}
\end{formal}
\begin{formal}
\begin{definition}[多元多项式全体]\label{def:多元多项式全体}
    记$\mathbb{K}$上的$n$元多形式全体为\[\mathbb{K}\left[x_1,x_2,\cdots,x_n\right]=
    \left\{
f\left(x_1,x_2,\cdots,x_n\right)=\sum_{\left(i_1,i_2,\cdots,i_n\right)\in\mathbb{
    N
}^n}a_{i_1i_2\cdots i_n}x_1^{i_1}x_2^{i_2}\cdots x_n^{i_n}
    \right\}\]
    这是一个$\mathbb{K}$-代数,称为
    关于$x_1,x_2,\cdots,x_n$的$n$元多项式代数（环）.
\end{definition}
\end{formal}
\begin{red}
\begin{remark}
    一个自然的问题是,$\mathbb{K}\left[x\right]$的各种性质,是否可以推广到$\mathbb{K}\left[x_1,x_2,\cdots,x_n\right]$.
\end{remark}
\end{red}
\begin{formal}
\begin{definition}[多元多项式的运算定义]\label{def:多元多项式的运算定义}
    加法定义为
    \[
    f\left(x\right)+g\left(x\right)=\sum_{\left(i_1,i_2,\cdots,i_n\right)}\left(a_{i_1i_2\cdots i_n}+b_{i_1i_2\cdots i_n}\right)x_1^{i_1}x_2^{i_2}\cdots x_n^{i_n}\]
    $\forall k\in\mathbb{K}$,数乘定义为
    \[
    kf\left(x\right)=\sum_{\left(i_1,i_2,\cdots,i_n\right)}\left(ka_{i_1i_2\cdots i_n}\right)x_1^{i_1}x_2^{i_2}\cdots x_n^{i_n} \]
    定义乘法则利用分配律
    转换为单项式的乘积之和并合并同类项.
\end{definition}
\end{formal}
\subsection{$n$元多项式代数环}
\begin{formal}
\begin{definition}[字典序]\label{def:字典序}
    将未定元按自然足标排序,即
    \[
    x_1\succ 
    x_2\succ\cdots\succ x_n
    \]
    然后对单项式排序,
    定义
    $\displaystyle
    ax_1^{i_1}x_2^{i_2}\cdots x_n^{i_n}\succ bx_{1}^{j_1}x_2^{j_2}\cdots x_n^{j_n}\left(ab\neq 0\right)$当且仅当
    $\exists 1\leqslant k\leqslant n,\textup{s.t.}i_1=j_1,i_2=j_2,\cdots,i_{k-1}=j_{k-1},i_k>j_k.$
\end{definition}
\end{formal}
\begin{green}
\begin{corollary}[字典序带来的唯一性]\label{cor:字典序带来的唯一性}
    在字典序下,每个多项式有唯一的排序,称为字典排序.
\end{corollary}
\end{green}
\begin{red}
\begin{remark}
    上述字典序的定义并不唯一.
\end{remark}
\end{red}
\begin{red}
\begin{remark}
    需要注意的是,在上述字典序下,首项未必是次数最高的单项式,而是字典序最大的单项式；末项未必是次数最低的单项式,而是字典序最小的单项式.
\end{remark}
\end{red}
\begin{green}
\begin{lemma}[首项维持]\label{lem:首项维持}
    在字典序下,$f\cdot g$的首项是$f$的首项与$g$的首项的乘积.
\end{lemma}
\begin{Proof}
若$f=0$或$g=0$,显然.

设$f$的首项为$ax_1^{i_1}x_2^{i_2}\cdots x_n^{i_n}$,$g$的首项为$bx_1^{j_1}x_2^{j_2}\cdots x_n^{j_n}$,则任取
$f\left(x\right)$的一个单项$cx_1^{k_1}x_2^{k_2}\cdots x_n^{k_n}$和
$g\left(x\right)$的一个单项$dx_1^{l_1}x_2^{l_2}\cdots x_n^{l_n}$,则$\exists 1\leqslant r,s\leqslant n\textup{s.t.}
i_1=k_1,i_2=k_2,\cdots,i_{r-1}=k_{r-1},i_r>k_r\Longleftrightarrow
a\succ c$且$j_1=l_1,j_2=l_2,\cdots,j_{s-1}=l_{s-1},j_s>l_s
\Longleftrightarrow b\succ d
$（其中系数指代整个单项
）,于是不妨设$r\leqslant s$则
\[
i_1+j_1=k_1+l_1,i_2+j_2=k_2+l_2,
\cdots,i_{r-1}+j_{r-1}=k_{r-1}+l_{r-1},i_r+j_r>k_r+l_r
\]
即
\[
ab x_1^{i_1+j_1}x_2^{i_2+j_2}\cdots x_n^{i_n+j_n}\succ cd x_1^{k_1+l_1}x_2^{k_2+l_2}\cdots x_n^{k_n+l_n}
\]
同理可证$ab\succ cb,ab\succ ad$,考虑到不可合并,则一定是首项.
\end{Proof}
\end{green}
\begin{formal}
\begin{theorem}[整性]\label{thm:整性}
    设$0\neq f,0\neq g\in\mathbb{K}\left[
    x_1,x_2,\cdots,x_n
    \right]$,则$f\cdot g\neq 0$.
\end{theorem}
\begin{Proof}
考虑首项即得.
\end{Proof}
\end{formal}
\begin{red}
\begin{remark}
    由\cref{thm:整性}整性立刻得到,多元多项式环满足乘法消去律.
\end{remark}
\end{red}
\begin{green}
\begin{corollary}[多元多项式的乘法消去律]\label{cor:多元多项式的乘法消去律}
    设 $f,g,h\in\mathbb{K}\left[x_1,x_2,\cdots,x_n\right]$,且 $f\neq 0$,则
    \[  f\cdot g=f\cdot h\Longrightarrow g=h \]
\end{corollary}
\end{green}
\subsection{齐次分解}
字典排序的缺点在于排序与单项式次数没有很好的关系,因此引入齐次分解.
\begin{formal}
\begin{definition}[齐次多项式]\label{def:齐次多项式}
    若多项式$f$的所有单项式均为$m$次单项式,
    则称$f$为$m$次齐次多项式或$m$次型.
\end{definition}
\end{formal}
\begin{formal}
    \begin{theorem}[齐次分解]\label{thm:齐次分解}
        \,
        \begin{enumerate}[label = {\textup{(\arabic*)}}]
            \item 若$f$、$g$均是$m$次型且$f+g\neq 0$,则$f+g$也是$m$次型
            \item 设$m$次型$f\neq 0 $和$k$次型$g\neq 0$,则$f\cdot g\neq 0$且$f\cdot g$是$m+k$次型
            \item 设$f\in\mathbb{K}\left[x_1,x_2,\cdots,x_n\right]$,则一定可以通过限制定义域得到与\[
            f=f_d+f_{d-1}+ \cdots+f_1+f_0
            \]
            其中$d=\deg\,f\geqslant 0,f_d\neq 0$为$d$次型,$f_i=0$或为 $i$次型$\left(\forall 0\leqslant i\leqslant d-1\right)$
        \end{enumerate}
    \end{theorem}
    \begin{Proof}
        \,
        \begin{enumerate}[label = {\textup{(\arabic*)}}]
            \item 若$f+g\neq 0$,则$f+g$的首项为$f$的首项与$g$的首项之和,故为$m$次型.
            \item 若$f\cdot g\neq 0$,则$f\cdot g$的首项为$f$的首项与$g$的首项之积,故为$m+k$次型.
            \item 逐步考查$d$次单项式、$d-1$次单项式$\cdots$即可.\qedhere
        \end{enumerate}
    \end{Proof}
\end{formal}
\begin{green}
\begin{lemma}[多元多项式运算的次数]\label{lem:多元多项式运算的次数}
    设$f$、$g\in\mathbb{K}\left[x_1,x_2,\cdots,x_n\right]$,则
    \begin{enumerate}[label = {\textup{(\arabic*)}}]
        \item  $\deg\,\left(f\cdot g\right)=\deg\,f+\deg\,g$
        \item $\deg \,\left(
            f\pm g
        \right)\leqslant \max\left\{
            \deg\,f,\deg\,g\right\}$
    \end{enumerate}
\end{lemma}
\begin{Proof}
考虑\cref{thm:齐次分解}齐次分解立得.\qedhere
\end{Proof}
\end{green}
\subsection{多元多项式与多元函数}
正如前文,多元多项式$f\left(
    x_1,x_2,\cdots,x_n
\right)\in
\mathbb{K}\left[
    x_1,x_2,\cdots,x_n\right]
$作赋值后成为一个多元函数$\overline{f}\left(
    x_1,x_2,\cdots,x_n
\right)$
\begin{align*}
    f:\mathbb{K}^n&\longrightarrow \mathbb{K}\\
    \begin{pmatrix}
        a_1\\
        a_2\\
        \vdots\\
        a_n
    \end{pmatrix}
       &\longmapsto f\left(
        a_1,a_2,\cdots,a_n
    \right)
\end{align*}
\begin{red}
    \begin{remark}
        一个自然的问题是,$f$、$g\in\mathbb{K}
        \left[
            x_1,x_2,\cdots,x_n
        \right]$作为多元多项式相等与作为多元函数相等是否等价.即下述关系是否成立
        \[
        f\left(
            x_1,x_2,\cdots,x_n
        \right)=g\left(
            x_1,x_2,\cdots,x_n
        \right)\Longleftrightarrow \overline{f}\left(
            x_1,x_2,\cdots,x_n
        \right)=\overline{g}\left(
            x_1,x_2,\cdots,x_n
        \right)
        \]  
    \end{remark}
\end{red}
\begin{red}
\begin{remark}
    由于对多元多项式相等的定义足够强,于是作为多元多项式相等可推得作为多元函数相等.
\end{remark}
\end{red}
\begin{green}
\begin{lemma}[非零多项式一定是非零函数]\label{lem:非零多项式一定是非零函数}
    设$0\neq f\in\mathbb{K}\left[
        x_1,x_2,\cdots,x_n
    \right]$,则存在$a_1,a_2,\cdots,a_n\in\mathbb{K}\,\textup{s.t.}$
    \[  f\left(
        a_1,a_2,\cdots,a_n
    \right)\neq 0 \]
\end{lemma}
\begin{Proof}
考虑归纳法,$n=1$显然,假设$n-1$成立,那么
只把$x_n$视为未定元,那么
\[
f=b_m\left(
    x_1,x_2,\cdots,x_{n-1}
\right)x_n^m+\cdots+b_1\left(
    x_1,x_2,\cdots,x_{n-1}
\right)x_n+b_0\left(
    x_1,x_2,\cdots,x_{n-1}
\right)
\]
其中$b_m\left(
    x_1,x_2,\cdots,x_{n-1}
\right)\neq 0$,那么由归纳假设,$\exists a_1,a_2,\cdots,a_{n-1}\in\mathbb{K}\,\textup{s.t.}$
\[
b_m\left(
    a_1,a_2,\cdots,a_{n-1}
\right)x_n^m+\cdots+b_1\left(
    a_1,a_2,\cdots,a_{n-1}
\right)x_n+b_0\left(
    a_1,a_2,\cdots,a_{n-1}
\right)\neq 0
\]
是一个关于$x_n$的$m$次多项式,故有根.
\end{Proof}
\end{green}
\begin{green}
\begin{corollary}[充要条件]\label{cor:充要条件}
     设$f$、$g\in\mathbb{K}\left[
        x_1,x_2,\cdots,x_n
    \right]$,则它们作为多元多项式相等当且仅当它们作为多元函数相等.
\end{corollary}
\begin{Proof}
作为多元多项式相等推得作为多元函数相等显然,作为多元函数相等推得作为多元多项式相等考虑反证法.即设$f$、$g\in\mathbb{K}\left[
    x_1,x_2,\cdots,x_n
\right]$且$\overline{f}=\overline{g},f\neq g$则$f-g\neq 0$,由上述引理,$\exists a_1,a_2,\cdots,a_n\in\mathbb{K}\,\textup{s.t.}$
\[
\left(f-g\right)
\left(
    a_1,a_2,\cdots,a_n
\right)=0\Longrightarrow
f\left(
    a_1,a_2,\cdots,a_n
\right)=g\left(
    a_1,a_2,\cdots,a_n
\right)
\]
矛盾,证毕.
\end{Proof}
\end{green}
\begin{red}
\begin{remark}
    同样的,可以定义多元多项式的整除关系、公因式、最大公因式、公倍式、最小公倍式、可约与不可约、
    辗转相除法、带余除法等等.但是,并不是所有的性质都可以推广到多元多项式.
\end{remark}
\end{red}
\begin{brown}
\begin{example}
    考虑$f=x,g=y\in\mathbb{K}\left[x,y\right]$,容易有$\left(f,g\right)=1$且$\forall 
    u\left(x,y\right),v\left(x,y\right)\in\mathbb{K}\left[x,y\right]$
    \[
     fu\left(x,y\right) +gv\left(x,y\right)\neq 1
    \]
\end{example}
\end{brown}
\begin{formal}
\begin{theorem}[多元多项式的因式分解]\label{thm:多元多项式的因式分解}
    设$f\left(
        x_1,x_2,\cdots,x_n
    \right)\in\mathbb{K}\left[
        x_1,x_2,\cdots,x_n
    \right]$,则有如下分解\[
    f\left(
           x_1,x_2,\cdots,x_n  
    \right)=
    cP_1\left(
        x_1,x_2,\cdots,x_n  
    \right)^{m_1}P_2\left(
        x_1,x_2,\cdots,x_n
    \right)^{m_2}\cdots P_k\left(
        x_1,x_2,\cdots,x_n
    \right)^{m_k}
    \]
    其中
    $ 
    c\in\mathbb{K},P_i \left(
        x_1,x_2,\cdots,x_n
    \right)\left(
        \forall 1\leqslant i\leqslant k
    \right)$是$\mathbb{K}$上的 $n$ 元不可约多项式.并且这一分解在相伴定义下是唯一的.
\end{theorem}
\end{formal}
\begin{red}
\begin{remark}
    \cref{thm:多元多项式的因式分解}的证明超出了高等代数的范围.
\end{remark}
\end{red}
\begin{red}
\begin{remark}
    基于\cref{thm:多元多项式的因式分解},我们将$n$ 元多项式环$\mathbb{K}\left[
        x_1,x_2,\cdots,x_n
    \right]$称为唯一分解整区（\textup{UFD}）.
\end{remark}
\end{red}
\newpage
\section{对称多项式}
\subsection{定义}
\begin{formal}
\begin{definition}[对称多项式]\label{def:对称多项式}
    设 $f\left(
        x_1,x_2,\cdots,x_n
    \right)\in\mathbb{K}\left[
        x_1,x_2,\cdots,x_n
    \right]$,若$\forall 1\leqslant i<j\leqslant n$,都有
    \[
    f\left(
        x_1,\cdots,x_i,\cdots,x_j,\cdots,x_n
    \right)=
    f\left(
        x_1,\cdots,x_j,\cdots,x_i,\cdots,x_n
    \right)
    \]则称$f\left(
        x_1,x_2,\cdots,x_n
    \right)$是$\mathbb{K}$上的$n$元对称多项式.
\end{definition}
\end{formal}
\begin{green}
\begin{lemma}[置换下定义]\label{lem:置换下定义}
    设$
     \left(
        k_1,k_2,\cdots,k_n
     \right)
    $是 $\left(
        1,2,\cdots,n
    \right)$的一个置换,且有对阵多项式
    $f\left(
        x_1,x_2,\cdots,x_n
    \right)$,则
    \[
    f\left(
        x_{k_1},x_{k_2},\cdots,x_{k_n}
    \right)=f\left(
        x_1,x_2,\cdots,x_n
    \right)
    \]
\end{lemma}
\begin{Proof}
根据\cref{thm:逆序数的实际意义},任意一个置换,经过其逆序数次相邻对换,就成为常排列,证毕.
\end{Proof}
\end{green}
\begin{green}
\begin{corollary}[对称多项式的和与积]\label{cor:对称多项式的和与积}
    对称多项式的和与积仍然是对称多项式.
\end{corollary}
\end{green}
\begin{green}
\begin{lemma}[对称多项式的复合]\label{lem:对称多项式的复合}
    设对称多项式$f_1,f_2,\cdots,f_m\in\mathbb{K}\left[
        x_1,x_2,\cdots,x_n
    \right]$,且有$g\left(
        y_1,y_2,\cdots,y_m
    \right)\in\mathbb{K}\left[
        y_1,y_2,\cdots,y_m
    \right]$,则$
    g\left(
        f_1,f_2,\cdots,f_m
    \right)
    $
    仍然是关于
    $x_1,x_2,\cdots,x_n$
    的对称多项式.
\end{lemma}
\end{green}
一个自然的问题是,我们能否用尽可能少的对称多项式去控制所有的对称多项式.
\subsection{初等对称多项式}
\begin{formal}
\begin{definition}[初等对称多项式]\label{def:初等对称多项式}
    设 $f\left(
        x_1,x_2,\cdots,x_n
    \right)\in\mathbb{K}\left[
        x_1,x_2,\cdots,x_n
    \right]$,则称
    \[
        \sigma_1\left(
            x_1,x_2,\cdots,x_n
        \right)=x_1+x_2+\cdots+x_n=
        \sum_{i=1}^{n}x_i\]\[
        \sigma_2\left(
            x_1,x_2,\cdots,x_n
        \right)=x_1x_2+x_1x_3+\cdots+x_{n-1}x_n=
        \sum_{1\leqslant i<\leqslant n}x_ix_j\]\[
        \cdots\cdots\cdots\cdots\cdots\cdots\]
        \[
        \sigma_r\left(
            x_1,x_2,\cdots,x_n
        \right)=
        \sum_{1\leqslant i_1<i_2<\cdots<i_r\leqslant n}x_{i_1}x_{i_2}\cdots x_{i_r}
        \]
        \[
        \cdots\cdots \cdots\cdots\cdots\cdots
        \]
        \[\sigma_n\left(
            x_1,x_2,\cdots,x_n
        \right)=x_1x_2\cdots x_n=
        \prod_{i=1}^{n}x_i
    \]
    为初等对称多项式.
\end{definition}
\end{formal}
\begin{formal}
\begin{theorem}[对称多项式基本定理]\label{thm:对称多项式基本定理}
    设$f\left(
        x_1,x_2,\cdots,x_n
    \right)$是$\mathbb{K}$上的$n$元对称多项式,则存在唯一的多项式
    $g\left(
        y_1,y_2,\cdots,y_n
    \right)$使得
    \[
    f\left(
        x_1,x_2,\cdots,x_n
    \right)=g\left(
        \sigma_1,\sigma_2,\cdots,\sigma_n
    \right)
    \]即任意一个对称多项式都可以唯一地表示为初等对称多项式的多项式.
\end{theorem}
\begin{Proof}
先证明存在性,设$f\neq 0$,则
取字典排序下首项为$ax_1^{i_1}x_2^{i_2}\cdots x_n^{i_n}\left(
    a\neq 0
\right)$,则断言$
i_1\geqslant i_2\geqslant\cdots\geqslant i_n
$（反证法,对换显然）.构造
\[
g_1\left(
    x_1,x_2,\cdots,x_n
\right)=
a\sigma_1^{i_1-i_2}\sigma_2^{i_2-i_3}\cdots\sigma_{n-1}^{i_{n-1}-i_n}\sigma_n^{i_n}
\]这是一个对称多项式,且其首项（
首项的乘积等于乘积的首项）为\[
ax_1^{i_1-i_2}\left(
    x_1x_2
\right)^{i_2-i_3}\cdots\left(
    x_1x_2\cdots x_n\right)^{i_n}=ax_1^{i_1}x_2^{i_2}\cdots x_n^{i_n}
\]

接下来做这样一件事,首先定义$f_0=f$,找出其首项$a
 x_1^{i_1}x_2^{i_2}\cdots x_n^{i_n}$,然后构造$g_1$,然后令$f_1=f_0-g_1$,则$f_1$的首项次数小于$f_0$,重复上述过程,对于一般的$f_i$,根据其首项
 $bx_1^{k_1} 
    x_2^{k_2}\cdots x_n^{k_n}\left(
        i_1\geqslant k_1\geqslant k_2\geqslant\cdots\geqslant k_n\geqslant 0
    \right)$构造$g_{i+1}$,
    $f_ {i+1}=f_i-g_{i+1}
     $
    这里的指数组$
    \left(
        k_1,k_2,\cdots,k_n
    \right)
    $
    的取法只有有限种.从而存在$s\in\mathbb{Z}^+\,\textup{s.t.}f_{s+1}=0.$反递推得到
    \[
    f=g_1+g_2+\cdots+g_s+g_{s+1}
    \]存在性证毕.

    再证明唯一性,设$h\left(
        y_1,y_2,\cdots,y_n
    \right)$使得
    \[
    f\left(
        x_1,x_2,\cdots,x_n
    \right)=h\left(
        \sigma_1,\sigma_2,\cdots,\sigma_n
    \right)
    \]令
    \[
    \varphi\left(
        y_1,y_2,\cdots,y_n  
    \right)=g\left(
        y_1,y_2,\cdots,y_n
    \right)-h\left(
        y_1,y_2,\cdots,y_n
    \right)
    \]则$
    \varphi\left(
        \sigma_1,\sigma_2,\cdots,\sigma_n   
    \right)=0
    $,只要证明$\varphi=0$即可.考虑反证法, $\varphi\neq 0$,因为
    \[
    \varphi =
    \cdots+c y_1^{j_1}y_2^{j_2}\cdots y_n^{j_n}+d y_1^{k_1}y_2^{k_2}\cdots y_n^{k_n}+\cdots
    \]
    无同类项且系数$\cdots,c,d,\cdots\neq 0$.则
    \begin{align*}
        0&=\varphi\left(
            \sigma_1,\sigma_2,\cdots,\sigma_n
        \right)\\
        &=\cdots+c\sigma_1^{j_1}\sigma_2^{j_2}\cdots\sigma_n^{j_n}+d\sigma_1^{k_1}\sigma_2^{k_2}\cdots\sigma_n^{k_n}+\cdots
    \end{align*}
    其中，因为$c\sigma_1^{j_1}
    \sigma_2^{j_2}\cdots\sigma_n^{j_n}
    $的首项是$
    cx_1^{j_1+j_2+\cdots+j_n}x_2^{j_2+\cdots+j_n}\cdots x_n^{j_n}
    $,而$d\sigma_1^{k_1}\sigma_2^{k_2}\cdots\sigma_n^{k_n}$的首项是$
    dx_1^{k_1+k_2+\cdots+k_n}x_2^{k_2+\cdots+k_n}\cdots x_n^{k_n}
    $.断言所有首项均不是同类项.考虑反证法,即有
    \[
    \begin{cases*}
        j_1+j_2+\cdots+j_n=k_1+k_2+\cdots+k_n\\
        j_2+\cdots+j_n=k_2+\cdots+k_n\\
        \qquad\cdots\cdots\cdots\cdots\\
        j_n=k_n
    \end{cases*}
    \]则$j_r=k_r,\forall 1\leqslant r\leqslant n$.既然所有的首项均不是同类项,那么一定可以找到一个字典排序下最先的首项,其先于$\varphi
    \left(
        \sigma_1,\sigma_2,\cdots,\sigma_n
    \right)$的所有其它单项,于是它不能够被消去,因此$
    \varphi\left(
        \sigma_1,\sigma_2,\cdots,\sigma_n
    \right)\neq 0
    $矛盾,证毕.
\end{Proof}
\end{formal}
\begin{brown}
\begin{example}
    设$f\left(x_1,x_2,x_3\right)=x_1^2x_2
    +x_1^2x_3
    +x_2^2x_1+x_2^2x_3+x_3^2x_1+x_3^2x_2$,求其初等对称多项式表示.
\end{example}
\begin{solution}
    显然,$f$的首项为$x_1^2x_2$,则$g_1=
    \sigma_1^{2-1}\sigma_2^{1-0}\sigma_3^{0}=f+3x_1x_2x_3=f+3\sigma_3
    \Longrightarrow f=\sigma_1\sigma_2-3\sigma_3.$
\end{solution}
\end{brown}
\begin{brown}
\begin{example}\label{ex:齐次对称多项式的初等对称多项式表示}
    设$6$次齐次对称
    多项式$f=\left(
        x_1^2+x_2^2
    \right)\left(
        x_2^2+x_3^2
    \right)\left(
        x_3^2+x_1^2
    \right)$,求其初等对称多项式表示.
\end{example}
\begin{solution}
显然$f$的首项为$x_1^4x_2^2$.进行上述构造时,某一$f_i$的首项
    $bx_1^{k_1}x_2^{k_2}x_3^{k_3}$满足：
    \[
    \begin{cases*}
    k_1+    k_2+k_3=6\\
    4\geqslant k_1\geqslant k_2\geqslant k_3\geqslant 0
    \end{cases*}
    \]
    于是该指数组$\left(
        k_1,k_2,k_3
    \right)$的有限取值及其$g_i$为
    \begin{align*}
       & \left(
        4,2,0
    \right)\Longrightarrow
    \sigma_1^{4-2}\sigma_2^{2-0}\sigma_3^{0}=
    \sigma_1^2
    \sigma_2^2,\\
    &\left(4,1,1\right) \Longrightarrow
    \sigma_1^{4-1}\sigma_2^{1-1}\sigma_3^{1}=
    \sigma_1^3\sigma_3,\\
    &\left(
        3,3,0   
    \right)
    \Longrightarrow
    \sigma_1^{3-3}\sigma_2^{3-0}\sigma_3^{0}=
    \sigma_2^3,\\
    &\left(
        3,2,1
    \right)
    \Longrightarrow
    \sigma_1^{3-2}\sigma_2^{2-1}\sigma_3^{1}=
    \sigma_1\sigma_2\sigma_3,\\
    &\left(
        2,2,2
    \right) \Longrightarrow
    \sigma_1^{2-2}\sigma_2^{2-2}\sigma_3^{2}=
    \sigma_3^2 
    \end{align*}
    且容易知道
    \[
    f=\sigma_1^2\sigma_2^2+a\sigma_1^3\sigma_3+b\sigma_2^3+c\sigma_1\sigma_2\sigma_3+d\sigma_3^2
    \]取特殊值得到
    \[
    a=-2,b=-2,c=4,d=-1
    \]
    于是
    \[
    f  =\sigma_1^2\sigma_2^2-2\sigma_1^3\sigma_3-2\sigma_2^3+4\sigma_1\sigma_2\sigma_3-\sigma_3^2
    \]
\end{solution}
\end{brown}
\begin{red}
\begin{remark}
    \cref{ex:齐次对称多项式的初等对称多项式表示}是齐次对称
    多项式的初等对称多项式表示,对于一般的对称多项式,可以进行齐次分解（容易证明分解得到的多项式均是对称的）后按\cref{ex:齐次对称多项式的初等对称多项式表示}方法进行.
\end{remark}
\end{red}
\subsection{Newton公式}
\begin{formal}
\begin{definition}[Fermal和]\label{def:Fermat和}
    定义
    \[
    s_k\left(
        x_1,x_2,\cdots,x_n
    \right)=x_1^k+x_2^k+\cdots+x_n^k\left(
        k\geqslant 1
    \right)
    \]为\textup{Fermat}和.显然是一个对称多项式.

    约定$S_0=n.$
\end{definition}
\end{formal}
\begin{green}
\begin{lemma}[Fermat与形式导数]\label{lem:Fermat与形式导数}
    设\begin{align*}
        f\left(x\right)&=
    \left(x-x_1\right)\left(x-x_2\right)\cdots\left(x-x_n\right)\\
    &=x^n-\sigma_1x^{n-1} +\sigma_2x^{n-2}+\cdots+\left(-1\right)^{n-1}
    \sigma_{n-1}x+\left(-1\right)^n\sigma_n
    \end{align*}则$\forall k\geqslant 1$
    \[
     x^{k+1}f'\left(x\right)=\left(
        s_0x^k+s_1 x^{k-1}+\cdots+s_k
     \right)f\left(x\right)+g\left(x\right)\]
     其中$\deg\,g\left(x\right)<n.$
\end{lemma}
\begin{Proof}
显然,\begin{align*}
  f'\left(x\right)&=\sum_{i=1}^{n}\left(x-x_1\right)\cdots\widehat{
      \left(x-x_i\right)   
} \cdots\left(x-x_n\right)  \\
&=\sum_{i=1}^{n}\frac{
    f\left(x\right)
}{x-x_i}
\end{align*}
则
\begin{align*}
    x^{k+1}f'\left(x\right)&=\sum_{i=1}^{n}x^{k+1}\frac{
        f\left(x\right)
    }{x-x_i}\\
    &=\sum_{i=1}^{n}\frac{
        x^{k+1}-x_i^{k+1}
    }{x-x_i}f\left(x\right)+\sum_{i=1}^{n}\frac{x_i^{k+1}f\left(x\right)}{x-x_i}
\end{align*}
则定义\[
g\left(x\right)=\sum_{i=1}^{n}\frac{x_i^{k+1}f\left(x\right)}{x-x_i} 
\]$\deg\,g\left(x\right)<n$.于是上式等于
\begin{align*}
    &\sum_{i=1}^{n}
    \left(
        x^{k}+x_ix^{k-1}+\cdots+x_i^k
    \right)f\left(x\right)+g\left(x\right)\\
    =&
    \left(
        s_0x^k+s_1 x^{k-1}+\cdots+s_k
    \right)f\left(x\right)+g\left(x\right)\qedhere
\end{align*}
\end{Proof}
\end{green}
\begin{formal}
\begin{theorem}[Newton公式]\label{thm:Newton公式}
记号同\cref{lem:Fermat与形式导数}
\begin{enumerate}  [label = {\textup{(\arabic*)}}]
    \item $k\leqslant n-1$,则
    \[
    s_k-s_{k-1}\sigma_1+\cdots+\left(-1\right)^{k-1}s_1\sigma_{k-1}+\left(-1\right)^k k\sigma_k=0\]
    \item $k\geqslant n$,则
    \[
    s_k-s_{k-1}\sigma_1+\cdots+
    \left(-1\right)^{n-1}s_{k - n+1}\sigma_{n-1}+
    \left(-1\right)^{n}s_{k-n}\sigma_n=0
    \]
\end{enumerate}
\end{theorem}
\begin{Proof}
因为
\[
f'\left(x\right)=
nx^{n-1}-\sigma_1\left(
    n-1
\right)x^{n-2}+\cdots+\left(-1\right)^{n-1}\sigma_{n-1}
\]
则
\[
x^{k+1} f'\left(x\right)=
nx^{n+k} -\sigma_1\left(
    n-1
\right)x^{n+k-1}+\cdots+\left(-1\right)^{n-1}\sigma_{n-1}x^{k+1}
\]
其中第$i$项是
$
\left(-1\right)^{n+k-i}\left(i-k\right)\sigma_{n+k-i} x^{i}
$,又因为
\[
x^{k+1}  f'\left(x\right)=
\left(
    s_0x^k+s_1x^{k-1}+\cdots+s_k
\right)\left(
    x^n-\sigma_1x^{n-1}+\cdots+\left(-1\right)^n\sigma_n
\right)+g\left(x\right)
\]接下来比较$x^n$的系数,$k\leqslant n-1$时为
\[
\left(-1\right)^k\left(n-k\right)\sigma_k
=s_k-s_{k-1}\sigma_1+\cdots+\left(-1\right)^{k}s_0\sigma_{k}
\]
其中$s_0=n$,于是
\[
s_k-s_{k-1}\sigma_1+\cdots+\left(-1\right)^{k-1}s_1\sigma_{k-1}+\left(-1\right)^k k\sigma_k=0
\]
$k\geqslant n$时,$x^n$的系数为
\[
0= s_k-s_{k-1}\sigma_1+\cdots+
\left(-1\right)^{n-1}s_{k - n+1}\sigma_{n-1}+
\left(-1\right)^{n}s_{k-n}\sigma_n\qedhere
\]
\end{Proof}
\end{formal}
\newpage
\section{结式与判别式}
\subsection{结式}
\begin{formal}
\begin{theorem}[基域扩张]\label{thm:基域扩张}
    设$f\left(x\right),g\left(x\right)\in\mathbb{K}\left[x\right]$,基于基域扩张下的不变性,有：$\left(
        f\left(x\right),g\left(x\right)
    \right)_{\mathbb{K}}=1\Longleftrightarrow
    \left(
        f\left(x\right),g\left(x\right)   
    \right)_{\mathbb{C}}=1\Longleftrightarrow
    f\left(x\right),g\left(x\right)$在$\mathbb{C}$无公共根.$\left(
        f\left(x\right),g\left(x\right)
    \right)\neq 1\Longleftrightarrow
    f\left(x\right),g\left(x\right)$在$\mathbb{C}$有公共根.
\end{theorem}
\end{formal}
\begin{green}
\begin{lemma}[公因子的存在性]\label{lem:公因子的存在性}
    设$f\left(x\right),g\left(x\right)\in\mathbb{K}\left[x\right],d\left(x\right)=\left(
        f\left(x\right),g\left(x\right)
    \right)$,则$d\left(x\right)\neq 1$当且仅当$\exists u\left(x\right)\neq 0,v\left(x\right)\neq 0\in\mathbb{K}\left[x\right]\,\textup{s.t.}$
    \[
    f\left(x\right)u\left(x\right)=g\left(x\right)v\left(x\right)
    \]其中
    $\deg \,u\left(x\right)<\deg\,g\left(x\right),\deg\,v\left(x\right)<\deg\,f\left(x\right)$.
\end{lemma}
\begin{Proof}
先考虑必要性,即设$d\left(x\right)\neq 1$,则
$f\left(x\right)=d\left(x\right)v\left(x\right),g\left(x\right)=d\left(x\right)u\left(x\right)$,显然成立.

再证充分性,即设$f\left(x\right)u\left(x\right)=g\left(x\right)v\left(x\right)$,则考虑反证法,设$\left(
    f\left(x\right),g\left(x\right)
\right)=1$,由
\[
f\left(x\right)\mid f\left(x\right)u\left(x\right)=g\left(x\right)v\left(x\right)
\]知$f\left(x\right)\mid v\left(x\right)$,这与
$v\left(x\right)\neq 0$且$\deg\,v\left(x\right)<
\deg\,f\left(x\right)$矛盾,证毕.
\end{Proof}
\end{green}
\begin{red}
\begin{remark}
    下设
    \begin{align*}
        &f\left(x\right)=a_0x^n +a_1x^{n-1}+\cdots+a_{n-1}x+a_n,a_0\neq 0\\
        &g\left(x\right)=b_0x^m+b_1x^{m-1}+\cdots+b_{m-1}x+b_m,b_0\neq 0\\
        &u\left(x\right)=x_0x^{m-1}+x_1x^{m-2}+\cdots+x_{m-2}x+x_{m-1}\\
        &v\left(x\right)=y_0x^{n-1}+y_1x^{n-2}+\cdots+y_{n-2}x+y_{n-1}
    \end{align*}
\end{remark}
\end{red}
代入\cref{lem:公因子的存在性}中的\[f\left(x\right)u\left(x\right)=g\left(x\right)v\left(x\right)\]得到一个包含$x_0,x_1,\cdots,x_{m-1},y_0,y_1,\cdots,y_{n-1}$共$m+n$个未知数的线性方程组
\[
\begin{cases*}
a_0x_0=  b_0y_0\\
a_1x_0+a_0x_1=b_1y_0+b_0y_1\\
a_2x_0+a_1x_1+a_0x_2=b_2y_0+b_1y_1+b_0y_2\\
\qquad \qquad   \qquad\qquad\cdots\cdots\cdots\cdots\\
a_nx_{m-3}+a_{n-1}x_{m-2}+a_{n-2}x_{m-1}=
b_m y_{n-3}+b_{m-1}y_{n-2}+b_{m-2}y_{n-1}\\
a_nx_{m-2}+a_{n-1}x_{m-1}=b_m y_{n-2}+b_{m-1}y_{n-1}\\
a_nx_{m-1}=b_m y_{n-1}
\end{cases*}
\]
其系数矩阵的转置
\[
\bm{A}'=\begin{pmatrix}
    a_0&a_1&a_2&\cdots&\cdots&a_n&0&\cdots&0\\
    0&a_0&a_1&\cdots&\cdots&a_{n-1}&a_n&\cdots&0\\
    0&0&a_0&\cdots&\cdots&a_{n-2}&a_{n-1}&\cdots&0\\
    \vdots&\vdots&\vdots&\vdots&\vdots&\vdots&\vdots&\vdots&\vdots\\
    0&0&\cdots&0&a_0&\cdots&\cdots&\cdots&a_n\\
    -b_0&-b_1&-b_2&\cdots&\cdots&\cdots&-b_{m}&\cdots&0\\
    0&-b_0&-b_1&\cdots&\cdots&\cdots&-b_{m-1}&-b_m&0\\
    \vdots&\vdots&\vdots&\vdots&\vdots&\vdots&\vdots&\vdots&\vdots\\
    0&\cdots&0&-b_0&-b_1&\cdots&\cdots&\cdots&-b_m
\end{pmatrix}
\]
因为$\left|\bm{A}'\right|=0$等价于该线性方程组存在非零解等价于
$u\left(x\right)\neq 0,v\left(x\right)\neq 0\in\mathbb{K}\left[x\right]$存在即等价于$f\left(x\right),g\left(x\right)$不互素.另一方面,若$\left|\bm{A}'\right|\neq 0$,则该线性方程组只有零解,即$f\left(x\right),g\left(x\right)$互素.
\begin{formal}
\begin{definition}[结式/Sylvester行列式]\label{def:结式}
    设
    \begin{align*}
        &f\left(x\right)=a_0x^n +a_1x^{n-1}+\cdots+a_{n-1}x+a_n,a_0\neq 0\\
        &g\left(x\right)=b_0x^m+b_1x^{m-1}+\cdots+b_{m-1}x+b_m,b_0\neq 0
    \end{align*}\newpage
    则称下列$m+n$阶行列式
    \[
    \left|R\left(f\left(x\right),g\left(x\right)\right)\right|=
    \begin{vmatrix}
    a_0&a_1&a_2&\cdots&\cdots&a_n&0&\cdots&0\\
    0&a_0&a_1&\cdots&\cdots&a_{n-1}&a_n&\cdots&0\\
    0&0&a_0&\cdots&\cdots&a_{n-2}&a_{n-1}&\cdots&0\\
    \vdots&\vdots&\vdots&\vdots&\vdots&\vdots&\vdots&\vdots&\vdots\\
    0&0&\cdots&0&a_0&\cdots&\cdots&\cdots&a_n\\
    b_0&b_1&b_2&\cdots&\cdots&\cdots&b_{m}&\cdots&0\\
    0&b_0&b_1&\cdots&\cdots&\cdots&b_{m-1}&b_m&0\\
    \vdots&\vdots&\vdots&\vdots&\vdots&\vdots&\vdots&\vdots&\vdots\\
    0&\cdots&0&b_0&b_1&\cdots&\cdots&\cdots&b_m
\end{vmatrix}
    \]为$f\left(x\right)$与$g\left(x\right)$的结式（\textup{resultant}）或\textup{Sylvester}行列式.

    若设
    \begin{align*}
       & \bm{\alpha} = (a_0, a_1, \ldots, a_{m-1}, a_m, \overbrace{0, \ldots, 0}^{\mathclap{n-1 \text{个}}})\\
&\bm{\beta} = (b_0, b_1, \ldots, b_{m-1}, b_m, \underbrace{0, \ldots, 0}_{n-1 \text{个}})
    \end{align*}
    定义一个循环映射
    \begin{align*}
        r:\mathbb{K}_{n+m}&\longrightarrow\mathbb{K}_{n+m}\\
        \left(c_1,c_2,\cdots,c_{n+m-1},c_{n+m}\right)&\longmapsto\left(c_{n+m},c_1,\cdots,c_{n+m-2},c_{n+m-1}\right)
    \end{align*}于是该结式可以写成
    \[
    R\left(f\left(x\right),g\left(x\right)\right)=\det\,(\begin{bmatrix}
        \bm{\alpha}\\r\left(\bm{\alpha}\right)\\\vdots\\r^{m-1}\left(\bm{\alpha}\right)\\\bm{\beta}\\r\left(\bm{\beta}\right)\\\vdots\\r^{n-1}\left(\bm{\beta}\right)
    \end{bmatrix})
    \]
\end{definition}
\end{formal}
\begin{red}
\begin{remark}
    此处记号的确混乱,更多时候,$R\left(
        f\left(x\right),g\left(x\right)
    \right)$是一个行列式.
\end{remark}
\end{red}
\begin{formal}
\begin{theorem}[结式判定互素]\label{thm:结式判定互素}
    $f\left(x\right),g\left(x\right)$在$\mathbb{C}$有公根当且仅当
    $R\left(f\left(x\right),g\left(x\right)\right)=0$.
\end{theorem}
\end{formal}
\begin{green}
\begin{corollary}[互素判定]\label{cor:互素判定}由\cref{thm:互素判定},
    $f\left(x\right),g\left(x\right)$在$\mathbb{C}$无公根当且仅当$R\left(
        f\left(x\right),g\left(x\right)
    \right)\neq 0$.
\end{corollary}
\end{green}
\begin{green}
\begin{lemma}[结式的两个变形]\label{lem:结式的两个变形}
    设常数$\lambda$则
    \[
    \left|
        R\left(
            f\left(x\right),g\left(x\right)\left(x-\lambda\right)
        \right)
    \right|=
    \left(-1\right)^nf\left(\lambda\right)R\left(f\left(x\right),g\left(x\right)\right)
    \]及
    \[
    \left|R\left(
        f\left(x\right),x-\lambda
    \right)\right|=\left(-1\right)^nf\left(\lambda\right)
    \]
\end{lemma}
\begin{Proof}
证明只需要简单的行列式变换.
\end{Proof}
\end{green}
\begin{formal}
\begin{theorem}[结式的根表示]\label{thm:结式的根表示}
    设
    \begin{align*}
        &f\left(x\right)=a_0x^n +a_1x^{n-1}+\cdots+a_{n-1}x+a_n,a_0\neq 0\\
        &g\left(x\right)=b_0x^m+b_1x^{m-1}+\cdots+b_{m-1}x+b_m,b_0\neq 0
    \end{align*}
    且$f\left(x\right)$的根为$x_1,x_2,\cdots,x_n$,$g\left(x\right)$的根为$y_1,y_2,\cdots,y_m$,则
    \[\left|
    R\left(f\left(x\right),g\left(x\right)\right)\right|=
    a_0^mb_0^n\prod_{i=1}^{n}\prod_{j=1}^{m}\left(
        x_i-y_j
    \right)
    \]
\end{theorem}
\begin{Proof}
采用摄动法证明.

第一步,首一化.$f\left(x\right)=a_0f_1\left(x\right),g\left(x\right)=b_0g_1\left(x\right)$,那么
\[\left|
R\left(f\left(x\right),g\left(x\right)\right)\right|=
a_0^mb_0^nR\left(f_1\left(x\right),g_1\left(x\right)\right)
\]于是考虑首一多项式的情况
\[\left|
R\left(f_1\left(x\right),g_1\left(x\right)\right)\right|
=
 \prod_{i=1}^{n}\prod_{j=1}^{m}\left(
        x_i-y_j
    \right)
\]
即可.

第二步,设$f\left(x\right)=\left(x-x_1\right)
\left(x-x_2\right)\cdots\left(x-x_n\right),g\left(x\right)=\left(x-y_1\right)
\left(x-y_2\right)\cdots\left(x-y_m\right)$.

则一方面,若$\exists 1\leqslant i\leqslant n,1\leqslant j\leqslant m\,\textup{s.t.}x_i=y_j,LHS=0=RHS.$

另一方面,设$x_i\neq y_j,\forall 1\leqslant i\leqslant n,1\leqslant j\leqslant m$且$x_1,x_2,\cdots,x_n$互不相同,$y_1,y_2,\cdots,y_m$互不相同.那么
\begin{align*}
    &R\left(f\left(x\right),g\left(x\right)\right)\begin{bmatrix}
        x_1^{n+m-1}&\cdots&x_n^{n+m-1}&y_1^{n+m-1}&\cdots&y_m^{n+m-1}\\
        x_1^{n+m-2}&\cdots&x_n^{n+m-2}&y_1^{n+m-2}&\cdots&y_m^{n+m-2}\\
        \vdots&&\vdots&\vdots&&\vdots\\
        x_1&\cdots&x_n&y_1&\cdots&y_m\\
        1&\cdots&1&1&\cdots&1
    \end{bmatrix}\\
    =&
    \begin{bmatrix}
        0&\cdots&0&y_1^{m-1}f\left(y_1\right)&\cdots&y_m^{m-1}f\left(y_m\right)\\
        \vdots&&\vdots&\vdots&&\vdots\\
        0&\cdots&0&f\left(y_1\right)&\cdots&f\left(y_m\right)\\
        x_1^{n-1}g\left(x_1\right)&\cdots&x_n^{n-1}g\left(x_n\right)&0&\cdots&0\\
        \vdots&&\vdots&\vdots&&\vdots\\
        g\left(x_1\right)&\cdots&g\left(x_n\right)&0&\cdots&0
    \end{bmatrix}
\end{align*}
取行列式得
\[
LHS = \left|
    R \left(
        f_1\left(x\right),g_1\left(x\right)    
    \right)
\right|\prod_{1\leqslant i<k\leqslant n}\left(
    x_i-x_k
\right)\prod_{1\leqslant j<l\leqslant m}\left(
    y_j-y_l 
\right)
\prod_{i=1}^{n}\prod_{j=1}^{m}\left(
    x_i-y_j
\right)
\]
利用\cref{thm:Laplace定理}\textup{Laplace}定理得到
\[
RHS=\left(-1\right)^{mn}\prod_{i=1}^{n}
g\left(x_i\right)\prod_{j=1}^{m}f\left(y_j\right)\prod_{1\leqslant i<k\leqslant n}\left(
    x_i-x_k
\right)\prod_{1\leqslant j<l\leqslant m}\left(
    y_j-y_l 
\right)
\]
注意到
\begin{align*}
    \left(-1\right)^{mn}\prod_{j=1}^{m}f\left(y_j\right)&=\prod_{j=1}^{m}\left(
        \left(-1\right)^n\left(y_j-x_1\right)\left(y_j-x_2\right)\cdots\left(y_j-x_n\right)
    \right)\\
    &=\prod_{j=1}^{m}\prod_{i=1}^{n}\left(
        x_i-y_j
    \right)
\end{align*}
于是
\begin{align*}
    \left|
    R\left(f_1\left(x\right),g_1\left(x\right)\right)
    \right|&=\prod_{i=1}^{n}g\left(x_i\right)\\
    &=
    \prod_{j=1}^{m}\prod_{i=1}^{n}\left(
        x_i-y_j
    \right)
\end{align*}至此两函数无同根及无同根且无重根的情况均证毕.

第三步,摄动法证明两函数无同根但有重根的情况.设$\exists c_1,c_2,\cdots,c_n,d_1,d_2,\cdots,d_m\in\mathbb{C}\,\textup{s.t.}\forall 0<t\ll 1,x_1+c_1t,x_2+c_2t,\cdots,x_n+c_nt,y_1+d_1t,y_2+d_2t,\cdots,y_m+d_mt$两两互不相同.

构造
\begin{align*}
    &f_t\left(x\right)=\left(
        x-x_1-c_1t
    \right)\left(
        x-x_2-c_2t
    \right)\cdots\left(
        x-x_n-c_nt
    \right)\\
    &g_t\left(x\right)=\left(
        x-y_1-d_1t
    \right)\left(
        x-y_2-d_2t
    \right)\cdots\left(
        x-y_m-d_mt
    \right)    
\end{align*}
二者的系数均是$t$的多项式,特别地,$f_0\left(x\right)=f\left(x\right),g_0\left(x\right)=g\left(x\right)$.由上一步,$\forall 0<t\ll 1$
\[
\left|R\left(
    f_t\left(x\right),g_t\left(x\right)    
\right)\right|=\prod_{i=1}^{n}
\prod_{j=1}^{m}
\left(
    x_i+c_it-y_j-d_jt
\right)
\]
此为关于$t$的连续函数,令$t\to 0^+$,则
\[
R\left(f\left(x\right),g\left(x\right)\right)=\prod_{i=1}^{n}
\prod_{j=1}^{m}
\left(
    x_i-y_j
\right)\qedhere
\]
\end{Proof}
\end{formal}
\subsection{判别式}
\begin{formal}
\begin{definition}[判别式]\label{def:判别式}
    设$f\left(x\right)=a_0x^n
    +a_1x^{n-1}+\cdots+a_{n-1}x+a_n\left(
        a_0\neq 0
    \right)$,则称\[
    \varDelta\left(f\right):=
    \left(-1\right)^{
        \frac{
            n\left(n-1\right)
        }{2}
    }a_0^{-1}R\left(
        f,f'
    \right)
    \]
    为$f\left(x\right)$的判别式.
\end{definition}
\end{formal}
\begin{formal}
\begin{theorem}[判别式的根定义]\label{thm:判别式的根定义}
    设 $f\left(x\right)=a_0x^n
    +a_1x^{n-1}+\cdots+a_{n-1}x+a_n\left(
        a_0\neq 0
    \right)$的根为
    $x_1,x_2,\cdots,x_n$
    ,则其判别式为
    \[
    \varDelta \left(f\right)=
    a_0^{2n-2}\prod_{1\leqslant i<j\leqslant n}\left(
        x_i-x_j
    \right)^2
    \]
\end{theorem}
\begin{Proof}
因为
\begin{align*}
R\left(
    f,g
\right)&=
a_0^m\prod_{i=1}^{n}\left(
    b_0\left(x_i-y_1\right)\left(x_i-y_2\right)\cdots\left(x_i-y_m\right)
\right)\\
&=a_0^m\prod_{i=1}^{n}g\left(x_i\right)
\end{align*}
令$g\left(x\right)=f'\left(x\right)\Longrightarrow m=n-1$,则
\[
f'\left(x_i\right)=a_0\left(x_i-x_1\right)\cdots\left(x_i-x_{i-1}\right)\left(x_i-x_{i+1}\right)\cdots\left(x_i-x_n\right)
\]于是
\begin{align*}
    R\left(f,f'\right)&=a_0^{n-1}\prod_{i=1}^{n}\left(a_0
        \left(x_i-x_1\right)\cdots\left(x_i-x_{i-1}\right)\left(x_i-x_{i+1}\right)\cdots\left(x_i-x_n\right)
    \right)\\
    &=\left(-1\right)^{
    \frac{n\left(n-1\right)}{2}
    }a_0^{2n-1}\prod_{1\leqslant i<j\leqslant n}\left(
        x_i-x_j
    \right)^2
\end{align*}
代入\cref{def:判别式}即得.
\end{Proof}
\end{formal}
\begin{green}
\begin{corollary}[重根判定]\label{cor:重根判定}
    多项式$f\left(x\right)$在$\mathbb{C}$有重根当且仅当其判别式$\varDelta\left(f\right)=0$.
\end{corollary}
\end{green}
\begin{red}
\begin{remark}
    之所以结式出现在多元多项式之后,是因为结式可以用于求解多元高次方程组.
\end{remark}
\end{red}
\begin{brown}
\begin{example}
    考虑求解
    \[
    \begin{cases*}
        f\left(x,y\right)=0\\
        g\left(x,y\right)=0
    \end{cases*}
    \]
\end{example}
\begin{solution}
设\begin{align*}
    &f(x, y) = a_0(y)x^n + a_1(y)x^{n-1} + \cdots + a_n(y)\\
&g(x, y) = b_0(y)x^m + b_1(y)x^{m-1} + \cdots + b_m(y)
\end{align*}
其中$a_i\left(y\right),b_j\left(y\right)$是$y$的多项式且$a_0\left(y\right)\neq 0,b_0\left(y\right)\neq 0.$
并设一解为$\left(
    \alpha,\beta
\right)$.令$y=\beta$,则$f\left(x,\beta\right)=0$与$g\left(x,\beta\right)=0$有公共根$x=\alpha$,则二者结式为零即
\newpage
\begin{align*}
    &R_x(f\left(x,\beta\right), g\left(x,\beta\right))\\ =&
\begin{vmatrix}
a_0(y) & a_1(y) & a_2(y) & \cdots&\cdots & a_n(y) & 0 & \cdots & 0 \\
0 & a_0(y) & a_1(y) & \cdots&\cdots & a_{n-1}(y) & a_n(y) & \cdots & 0 \\
0 & 0 & a_0(y) & \cdots&\cdots & a_{n-2}(y) & a_{n-1}(y) & \cdots & 0 \\
\vdots & \vdots & \vdots & \vdots&\vdots & \vdots & \vdots & \vdots & \vdots \\
0 & 0 & \cdots & 0 & a_0(y) & \cdots&\cdots&\cdots & a_n(y) \\
b_0(y) & b_1(y) & b_2(y) & \cdots&\cdots&\cdots & b_m(y) &  \cdots & 0 \\
0 & b_0(y) & b_1(y) & \cdots&\cdots&\cdots & b_{m-1}(y) & b_m(y) & \cdots \\
0 & 0 & b_0(y) & \cdots & b_{m-2}(y) & b_{m-1}(y) & \cdots & b_m(y) \\
\vdots & \vdots & \vdots & \vdots & \vdots & \vdots & \vdots&\vdots & \vdots \\
0 & \cdots & 0 & b_0(y) & b_1(y) & \cdots&\cdots&\cdots & b_m(y)
\end{vmatrix}\\
= &0
\end{align*}
这是一个关于$y$的多项式,一定有根$y=\beta$,反代回去就完成了未定元的减少.
\end{solution}
\end{brown}
\begin{red}
\begin{remark}
    从上面可以看出一般的思路,即先求$R_x\left(f,g\right)$（一个关于$y$的多项式）,再求其根$\beta_1,\cdots,\beta_r$,再代入$f\left(x,\beta_i\right)=0$解得$x=\alpha_1,\cdots,\alpha_s$,最后,需要验证$\left(
        \alpha_i,\beta_j
    \right)$是否是合理解.
\end{remark}
\end{red}